\documentclass[11pt]{article}

\usepackage[margin=1in]{geometry}
\usepackage{amsmath,amssymb,amsthm,mathtools}
\usepackage{microtype}
\usepackage[round]{natbib}
\usepackage{float}
\usepackage{xcolor}
\usepackage{todonotes}
\usepackage{aliascnt}
\usepackage[colorlinks=true,linkcolor=blue!55!black,citecolor=blue!55!black,urlcolor=blue!55!black,hypertexnames=false]{hyperref}
\usepackage[capitalise,noabbrev]{cleveref}
\usepackage{autonum} % load last

\newtheorem{theorem}{Theorem}
\newaliascnt{proposition}{theorem}
\newtheorem{proposition}[proposition]{Proposition}
\aliascntresetthe{proposition}
\newaliascnt{lemma}{theorem}
\newtheorem{lemma}[lemma]{Lemma}
\aliascntresetthe{lemma}
\newaliascnt{corollary}{theorem}
\newtheorem{corollary}[corollary]{Corollary}
\aliascntresetthe{corollary}

\crefname{equality}{equality}{equalities}
\Crefname{equality}{Equality}{Equalities}
\crefname{inequality}{inequality}{inequalities}
\Crefname{inequality}{Inequality}{Inequalities}
\newcommand{\eqlabel}[1]{\label[equality]{#1}}
\newcommand{\ineqlabel}[1]{\label[inequality]{#1}}
\newcommand{\proofstep}[2]{%
    \par\medskip
    \noindent
    \textcolor{blue!55!black}{\rule{2pt}{1.2em}}%
    \hspace{0.6em}%
    {\sffamily\bfseries\color{blue!55!black}Step #1}%
    \quad
    {\sffamily #2}%
    \par\smallskip
}
\newcommand{\proofstepinline}[1]{%
    {\sffamily\bfseries\color{blue!55!black}Step #1.}\enspace
}
\newcommand{\R}{\mathbb R}
\newcommand{\E}{\mathbb E}
\newcommand{\Pp}{\mathbb P}
\newcommand{\K}{\mathcal K}
\newcommand{\cP}{\mathcal P}
\newcommand{\F}{\mathcal F}
\newcommand{\N}{\mathcal N}
\newcommand{\T}{\mathcal T}
\newcommand{\tr}{\operatorname{tr}}
\newcommand{\Lip}{\operatorname{Lip}}
\newcommand{\op}{\mathrm{op}}
\newcommand{\argmin}{\operatornamewithlimits{arg\,min}}
\newcommand{\1}{\mathbf 1}
\newcommand{\eps}{\varepsilon}
\renewcommand{\epsilon}{\varepsilon}
\newcommand{\ip}[2]{\left\langle #1,#2\right\rangle}
\newcommand{\norm}[1]{\left\lVert #1\right\rVert}
\newcommand{\normm}[2]{\left\lVert #1\right\rVert_{#2}}
\newcommand{\alireza}[1]{\todo[color=blue!20,size=\tiny]{Alireza: #1}}

\floatstyle{ruled}
\newfloat{algorithm}{tbp}{loa}
\floatname{algorithm}{Algorithm}

\title{Direct Hermite Projections for Stochastic Bandit Convex Optimization}
\author{\qquad}
\date{}

\begin{document}
\maketitle

\begin{abstract}
    We consider stationary stochastic bandit convex optimization with a bounded
    convex loss and subgaussian observation noise.  Gaussian online Newton step
    methods are usually analyzed through an optimistic Gaussian surrogate whose
    derivatives are estimated using a likelihood-ratio correction.  We show that
    neither device is needed in the stationary problem.  The algorithm estimates
    the first two Gaussian Hermite coefficients of a queryable convex extension
    directly from the observed loss and applies online Newton step to the resulting
    (damped) quadratic estimates.
    For a rounded convex body $\K$, the regret is
    \begin{align}
        O\!\left(d^2 M\sqrt n\,
        \log^4\!\frac{nd}{\delta}\right)
    \end{align}
    with probability at least $1-\delta$, where $M$ is $\tilde{O}(d^{-1/2})$
    when $\frac1d \K$ is in L\"owner position,
    which recovers the stationary rate of \citet{lattimore2023second}
    of $\widetilde O(d^{3/2}\sqrt n)$.
    % In this case, the regret is $\widetilde O(d^{3/2}\sqrt n)$.
\end{abstract}

\section{Introduction}

The learner in bandit convex optimization chooses a point in a convex body
and observes only a noisy value of the loss at that point.  A useful way to
extract second-order information is to sample from a Gaussian and estimate the
first two coefficients in the corresponding Hermite expansion.  The first
coefficient behaves like a gradient and the second like an average Hessian, so
the natural update is online Newton step.

There is one difficulty with this direct idea.  The online Newton step
potential contains the curvature of the quadratic at the comparator, while
convexity controls regret through a Bregman divergence of the original loss.
These quantities need not be comparable globally.  The usual Gaussian ONS
analysis resolves this by first replacing the loss with an optimistic
Gaussian surrogate and then estimating its derivatives with a change of
measure; see \citet{lattimore2023second} and Chapters~10--12 of
\citet{lattimore2026bandit}.  Here we take a
different route for the stationary stochastic problem.  We compare the direct
Hermite curvature with the Bregman divergence on an ellipsoidal region around
the Gaussian mean.  The comparison loses only a small part of the usual
log-determinant trace budget.

The resulting algorithm has one Gaussian sample, one function evaluation and
one metric projection per round.
The proof still has to address two features of bandit feedback:
the one-sample Hessian estimate is indefinite, and the queryable extension
amplifies noise outside the constraint set.
We handle the first by comparing the random covariance with its
predictable counterpart, and the second by using the extra curvature of the
standard reinforced extension.

The paper is organised as follows.  \Cref{sec:setting} introduces the problem,
the extension and the algorithm.
Trace, moment and concentration estimates
are collected in \cref{sec:stochastic-control}.  The stability statement and
regret argument appear in \cref{sec:stability}, with the supporting proofs
deferred to the appendices.

\section{Problem formulation and algorithm}
\label{sec:setting}

For $r>0$, let $\mathbb{B}_r=\{x\in\R^d:\norm{x}\leq r\}$
and $\mathbb{S}_r=\{x\in\R^d:\norm{x}=r\}$ be the Euclidean ball and sphere of radius $r$.
Let $\K\subset\R^d$ be a convex body.  An affine preprocessing step puts the body
in the position $\mathbb{B}_1\subseteq\K\subseteq \mathbb{B}_{2d}$.
The loss $f:\K\to[0,1]$ is convex and does not change with time.  In round
$t$, the learner chooses $A_t\in\K$ and observes $f(A_t)+\xi_t$
where $\xi_t$ is the noise in round $t$.
Let $\zeta_t$ be the randomness of the algorithm in round $t$ which it uses to make decisions, and let
$\F_t=\sigma(\zeta_1,\xi_1,\ldots,\zeta_t,\xi_t)$ be the filtration at the end of round $t$,
and $\E_t[\cdot]=\E[\cdot\mid\F_{t-1}]$ and $\Pp_t(\cdot)=\Pp(\cdot\mid\F_{t-1})$ be the conditional expectation and probability given the history.
The noise is conditionally subgaussian, i.e., for all $t$,
\begin{align}
    \E_t[\xi_t|\zeta_t]=0,
    \qquad
    \E_t[\exp(\xi_t^2)|\zeta_t]\leq2.
    \label{eq:noise}
\end{align}
The quantity of interest is the realized (\textit{pseudo-})regret
\begin{align}
    \operatorname{Reg}_n
    =\sum_{t=1}^n\left(f(A_t)-\inf_{x\in\K}f(x)\right).
    \label{eq:regret}
\end{align}

Let $\pi(x)=\inf\{r>0:x\in r\K\}$ be the Minkowski functional of $\K$ and let
$U$ be uniform on the Euclidean unit sphere.
The only geometric quantity in the result is
\begin{align}
    M=\max\left\{\frac1{\sqrt d},\E[\pi(U)]\right\}.
    \label{eq:M}
\end{align}
The rounding  $\mathbb{B}_1\subseteq\K$ implies that $\pi$ is $1$-Lipschitz.
Write $\gamma_{\mu,\Sigma}$ for the density of
$\N(\mu,\Sigma)$, for $\mu\in\R^d$ and $\Sigma\in\R^{d\times d}$ symmetric positive semidefinite.

\subsection{Distributional derivatives}
A general convex function $h:\R^d\to\R$ may not be differentiable everywhere and
its Hessian may not exist in the classical sense.
Traditionally, this issue is handled through a mollification process: convolve $h$ with a smooth kernel to obtain a smooth approximation, and then take limits.
However, we will take a different approach that avoids mollification and works directly with the original function $h$.
We employ the theory of distributions, which allows us to define derivatives in a weak sense.
Practically, first we recall the definition of the Schwartz space $\mathcal S(\R^d)$, which consists of smooth functions that decay faster than any polynomial at infinity, along with all their derivatives.
We write $Dh$ and $D^2h$ for the first and second distributional derivatives of $h$, which
are defined as linear functionals on the Schwartz space $\mathcal S(\R^d)$ of smooth rapidly decaying functions.
\begin{align}
    \langle Dh,\varphi\rangle
     & =-\int_{\R^d}h(x)\nabla\varphi(x)\,dx,
     &
    \langle D^2h,\varphi\rangle
     & =\int_{\R^d}h(x)\nabla^2\varphi(x)\,dx,
    \qquad
    \varphi\in\mathcal S(\R^d),
\end{align}
where $\varphi$ is the Gaussian density $\gamma_{\mu,\Sigma}$ for our purposes,
which is a Schwartz function.
The crucial fact here is that if $h$ is convex, then $D^2h$ is a positive-semidefinite matrix-valued \textit{tempered} Radon measure and the Gaussian density $\gamma_{\mu,\Sigma}$ is a Schwartz function, so the pairing $\langle D^2h,\gamma_{\mu,\Sigma}\rangle$ is well-defined and yields a positive-semidefinite matrix.
\alireza{This subsection needs more explanations and citations.}

\subsection{A queryable convex extension}

Fix $\eps\in(0,1/2]$ and write $\K_\eps=(1-\eps)\K$ and
\begin{align}
    \pi_\eps(x)
     & =\max\left\{1,\frac{\pi(x)}{1-\eps}\right\}
    =\begin{cases}
         1               & \text{if }x\in\K_\eps, \\
         \pi(x)/(1-\eps) & \text{otherwise}\,,
     \end{cases}
    \quad\text{and}\quad v(x)=\pi_\eps(x)-1\,.
    \label{eq:extension-notation}
\end{align}
We use the queryable extension from
\citet[Eq.~(10.6)]{lattimore2026bandit} which is defined for all $x\in\R^d$ by
\begin{align}
    e(x)
    =\pi_\eps(x)f\!\left(\frac{x}{\pi_\eps(x)}\right)
    +\frac{2}{\eps}v(x).
    \label{eq:extension}
\end{align}
\begin{lemma}[{\citet[Proposition~3.23]{lattimore2026bandit}}]
    \label{lem:extension}
    The function $\bar{e}:
        x
        \longmapsto
        \pi_\eps(x)
        f\!\left(\frac{x}{\pi_\eps(x)}\right)
        +
        \frac{1-\eps}{\eps}v(x)
    $ is convex, agrees with $f$ on $\K_\eps$, and
    \begin{align}
        \Lip(\bar e)\leq\frac{2}{\eps(1-\eps)}\,.
    \end{align}
    Moreover,
    $e\!\left(\frac{x}{\pi_\eps(x)}\right)\leq e(x)$
    for all $x\in\R^d$.
\end{lemma}
Now, observe that $e(x)=\bar e(x)+\frac{1+\eps}{\eps}v(x)$, so $e$ is also convex, agrees with $f$ on $\K_\eps$ since $v(x)=0$ for $x\in\K_\eps$.
Moreover, $v$ is $1/(1-\eps)$-Lipschitz, so $e$ is also Lipschitz with
\begin{align}
    \Lip(e)\leq\frac{2}{\eps(1-\eps)}+\frac{1+\eps}{\eps}\cdot\frac{1}{1-\eps}\leq\frac{8}{\eps}\,,
    \label{eq:extension-lipschitz}
\end{align}
when $\eps\in(0,1/2]$.

\subsection{Gaussian $L^2$ projection and Hermite basis}
\label{sec:hermite-projection}
The Gaussian space $L^2(\N(\mu,\Sigma))$ consists of the functions that are
square-integrable under the law of $X \sim \N(\mu, \Sigma)$, i.e.
$L^2(\N(\mu,\Sigma))
    =
    \left\{
    r:\R^d\to\R:
    \E[r(X)^2]<\infty
    \right\}$.
Since $e$ is Lipschitz, we have
\begin{align}
    \E\left[
        e(X)^2
        \right]
    \leq
    \E\left[
        (e(0)+\Lip(e)\norm{X})^2
        \right]
    \leq
    2e(0)^2 + 2\Lip(e)^2\E[\norm{X}^2] < \infty,
\end{align}
which implies that $e\in L^2(\N(\mu,\Sigma))$.
Let $\mathcal Q_2$ be the space of quadratic polynomials on $\R^d$,
each of which can be written uniquely as
\begin{align}
    x\mapsto c+\ip{a}{x}+\frac12 x^\top Bx,
    \qquad
    c\in\R,\quad a\in\R^d,\quad B\in\mathbb{R}_{\text{sym}}^{d\times d}.
    \label{eq:quadratic-polynomial}
\end{align}
It is easy to check that $\mathcal Q_2\subset L^2(\N(\mu,\Sigma))$.
Through defining the inner product $\ip{r}{s}_{L^2(\N(\mu,\Sigma))}=\E[r(X)s(X)]$, the space $L^2(\N(\mu,\Sigma))$ becomes a Hilbert space, and we can define the projection of $e$ onto the space of quadratic polynomials as
\begin{align}
    \mathcal P_{\mu,\Sigma}e
    =
    \argmin_{p\in\mathcal Q_2}
    \E\left[
        (e(X)-p(X))^2
        \right]\,.
    \label{eq:hermite-best-approximation}
\end{align}
\begin{lemma}
    \label{lem:hermite-projection-convex}
    Let $e:\R^d\to\R$ be a Lipschitz convex function.
    Then its projection $\mathcal P_{\mu,\Sigma}e$ is
    \begin{align}
        \mathcal P_{\mu,\Sigma}e(x)
        =
        \E[e(X)]
        +\ip{g}{x-\mu}
        +\frac12
        \left(
        \|x-\mu\|^2_{H}
        -\tr(H\Sigma)
        \right),
    \end{align}
    where
    $g = \langle De, \gamma_{\mu,\Sigma}\rangle$,
    and
    $H = \langle D^2e, \gamma_{\mu,\Sigma}\rangle$.
    \alireza{We might as well use $\langle e, D^2\gamma_{\mu,\Sigma}\rangle=\E[e(x)\nabla^2\gamma_{\mu,\Sigma}(x)]$
        as the definition of $H$ and avoid the distributional derivative in the main paper.}
\end{lemma}
\begin{proof}
    We need a proof, although we don't really need the \textit{projection} property.\alireza{
        Technically we can drop this lemma and just say that we define
        $\mathcal P_{\mu,\Sigma}e(x)$ as such without any justification.
        But it is nice to have a lemma that says that this is indeed the Hermite projection.
    }
    Lipschitzness of $e$ implies the temperedness of $e$, which ensures that $\langle De, \gamma_{\mu,\Sigma}\rangle$ and $\langle D^2e, \gamma_{\mu,\Sigma}\rangle$ are well-defined.
    Convexity of $e$ implies that $D^2e$ is a positive-semidefinite matrix-valued tempered Radon measure, which ensures that $H$ is positive-semidefinite and justifies the notation $\|x-\mu\|^2_{H} = (x-\mu)^\top H (x-\mu)$.
\end{proof}






\section{The algorithm}

On an active round $t$, the learner samples
$X_t\sim \N(\mu_t,\Sigma_t)$ and plays its radial contraction
$A_t=X_t/\pi_\eps(X_t)\in\K$.
The resulting feedback is $f(A_t) + \xi_t$ which is converted into the feedback from the queryable extension $e$ as
\begin{align}
    Y_t
     & =
    \pi_\eps(X_t)\bigl(f(A_t)+\xi_t\bigr)
    +\frac2\eps\bigl(\pi_\eps(X_t)-1\bigr)
    =e(X_t)+\pi_\eps(X_t)\xi_t,
    \qquad
    \E_t[Y_t\mid X_t]=e(X_t).
    \label{eq:Y}
\end{align}
Let $\lambda \in (0,1)$ be a damping parameter and define the quadratic surrogate
\alireza{
    The connection to the previous section is not clear.
    There is also the sudden presence of the damping parameter $\lambda$ in the definition of $q_t$,
    which should be explained. I think it is better to translate this quantity into somewhere else in the algorithm and the analysis to make the story more coherent.
}
\begin{align}
    q_t(x)
    =
    \ip{g_t}{x-\mu_t}
    +\frac\lambda2\normm{x-\mu_t}{H_t}^2,
    \qquad
    g_t\triangleq \langle De,\gamma_{\mu_t,\Sigma_t}\rangle,
    \qquad
    H_t\triangleq \langle D^2e,\gamma_{\mu_t,\Sigma_t}\rangle.
\end{align}
\alireza{
    Again, $g_t$ and $H_t$ can be defined through expectations to avoid distributional notation in the main paper.
}
The main observation is that $g_t$ and $H_t$ can be estimated from the single pair $(X_t, Y_t)$.
More precisely, define
\begin{align}
    \hat{q}_t(x)
    \triangleq
    \ip{\hat{g}_t}{x-\mu_t}
    +
    \frac\lambda2(x-\mu_t)^\top\hat{H}_t(x-\mu_t)\,,
\end{align}
where for $Z_t\triangleq\Sigma_t^{-1/2}(X_t-\mu_t)$, we define
\begin{align}
    \widehat g_t
     & \triangleq - Y_t \nabla \gamma_{\mu_t,\Sigma_t}(X_t)
    =
    Y_t\Sigma_t^{-1/2}Z_t,
     &
    \widehat H_t \triangleq Y_t \nabla^2 \gamma_{\mu_t,\Sigma_t}(X_t)
     & =Y_t\Sigma_t^{-1/2}(Z_tZ_t^\top-I_d)\Sigma_t^{-1/2}\,,
    \label{eq:sampled-hermite-coefficients}
\end{align}
then from distributional integration by parts\alireza{More explanation needed here.} we get
\begin{align}
    g_t = \langle De, \gamma_{\mu_t,\Sigma_t}\rangle
    = - \langle e, \nabla \gamma_{\mu_t,\Sigma_t}\rangle
    = \E_t[\widehat g_t],
    \quad\text{and}\quad
    H_t = \langle D^2e, \gamma_{\mu_t,\Sigma_t}\rangle
    = \langle e, \nabla^2 \gamma_{\mu_t,\Sigma_t}\rangle
    = \E_t[\widehat H_t]\,,
    \label{eq:sampled-hermite-unbiasedness}
\end{align}
and more importantly, $\E_t[\hat{q}_t\mid X_t]=q_t$.

The algorithm is just projected online newton steps (ONS) applied to quadratic estimates $\widehat q_t$.\alireza{explanation on ONS is needed.}
If $\Sigma_t^{-1}+\lambda\eta\widehat H_t\succ0$,
then completing the square gives the equivalent forms
\begin{align}
    \mu_{t+1}
     & =\argmin_{\mu\in\K_\eps}
    \left\{
    \frac12\normm{\mu-\mu_t}{\Sigma_t^{-1}}^2
    +\eta\widehat q_t(\mu)
    \right\}
    =\argmin_{\mu\in\K_\eps}
    \normm{\mu-(\mu_t-\eta\Sigma_{t+1}\widehat g_t)}
    {\Sigma_t^{-1}+\lambda\eta\widehat H_t}^2.
    \label{eq:projected-ons-update}
\end{align}
\alireza{It should be clear from the update why the bias term $\E_t[e(X_t)]$ is dropped from $q_t$
    as it does not affect the Newton step.}
\begin{algorithm}[!t]
    \caption{Direct Hermite--ONS algorithm}
    \label{alg:direct-hermite}
    \small
    \begin{tabbing}
        \qquad\=\qquad\=\kill
        \textbf{input:} $\eta,\lambda,\sigma^2,\eps$\\
        $\mu_1=0$, $\Sigma_1=\sigma^2I_d$ and $\mathsf{safe}=\mathsf{true}$\\[1mm]
        \textbf{for} $t=1,\ldots,n$ \textbf{do}\\
        \>Sample $X_t\sim\N(\mu_t,\Sigma_t)$\\
        \>Play $A_t=X_t/\pi_\eps(X_t)$ and observe $f(A_t)+\xi_t$\\
        \>Form $Y_t$, $\widehat g_t$, and $\widehat H_t$ as in
        \eqref{eq:Y} and \eqref{eq:sampled-hermite-coefficients}\\
        %  and set $V_{t+1}=\Sigma_t^{-1}+\lambda\eta\widehat H_t$\\
        \>\textbf{if} $\Sigma_t^{-1}+\lambda\eta\widehat H_t\succ0$:\\
        \>\>Set $\Sigma_{t+1}=(\Sigma_t^{-1}+\lambda\eta\widehat H_t)^{-1}$ and update
        $\mu_{t+1}$ as in \eqref{eq:projected-ons-update}\\
        \>\textbf{else}:\\
        \>\>Set $\Sigma_{t+1}=\Sigma_t$ and $\mu_{t+1}=\mu_{t}$ \qquad\textit{\(\triangleright\) The algorithm has failed.}\\
        \textbf{end for}
    \end{tabbing}
\end{algorithm}

\subsection{Main result}

Fix $\delta\in(0,1/2)$ and let
\begin{align}
    \ell=C_0\log\frac{C_0nd}{\delta},
    \eqlabel{eq:ell}
\end{align}
where $C_0$ is a sufficiently large universal constant.

\begin{theorem}
    \label{thm:main}
    There are universal constants $c_\lambda,c_\eta,C_\eps,C>0$ for which the
    following holds.  Set
    \begin{align}
        \sigma  & =\frac1{M\sqrt{2d}},
                &
        \lambda & =\frac{c_\lambda}{d^{3/2}M\ell^2},
        \label{eq:parameters-one}                     \\
        \eta    & =\frac{c_\eta dM}{\sqrt n},
                &
        \eps_0  & =\frac{C_\eps d^2M\ell^4}{\sqrt n}.
        \label{eq:parameters-two}
    \end{align}
    If $\eps_0\leq1/2$, run \cref{alg:direct-hermite} with $\eps=\eps_0$.
    Otherwise play any fixed point in $\K$.  In the setting above,
    \begin{align}
        \operatorname{Reg}_n
        \leq C d^2M\ell^4\sqrt n
        \label{eq:main-regret}
    \end{align}
    with probability at least $1-\delta$.
\end{theorem}

\begin{corollary}
    \label{cor:loewner}
    If $\K/d$ is in L\"owner position, then
    \begin{align}
        \operatorname{Reg}_n
        =O\!\left(
        d^{3/2}\sqrt{\log(ed)}\,\ell^4\sqrt n
        \right)
        =\widetilde O(d^{3/2}\sqrt n)
    \end{align}
    with probability at least $1-\delta$.
\end{corollary}

\begin{proof}
    In L\"owner position,
    $\E[\pi(U)]=O(\sqrt{\log(ed)/d})$; see
    \citet[Table~10.1(a)]{lattimore2026bandit}.  Substitute this bound in
    \eqref{eq:main-regret}.
\end{proof}

\section{Analysis}
\label{sec:analysis}
\subsection{Regret switching}
We fix a time step $t$ and drop the subscript $t$ from the notation in this subsection.
The proof rests on relating the regret of the quadratic surrogate $q$ to the regret of the original loss $e$.
For a convex Lipschitz function $e$, and a differentiable point $z$, the Bregman divergence is
\begin{align}
    D_e(x\,\|\,z)
    =e(x)-e(z)-\ip{\nabla e(z)}{x-z}.
    \label{eq:bregman}
\end{align}

\begin{lemma}
    \label{lem:regret-switch}
    Let $X\sim\N(\mu,\Sigma)$, then for every $x\in\R^d$,
    \begin{align}
        \E[e(X)]-e(x)
        =\tr(H\Sigma)+\ip{g}{\mu-x}-\E[D_e(x\,\|\,X)]\,,
        \label{eq:regret-switch}
    \end{align}
    where $g=\langle De,\gamma_{\mu,\Sigma}\rangle$ and
    $H=\langle D^2e,\gamma_{\mu,\Sigma}\rangle$.
\end{lemma}
\begin{proof}
    Splitting $X-x=(X-\mu)+(\mu-x)$ in a rearrangement of \eqref{eq:bregman} and taking expectations gives
    \begin{align}
        \E[e(X)]-e(x)
         & =
        \E\left[\ip{\nabla e(X)}{X-\mu}\right]
        +
        \ip{\E[\nabla e(X)]}{\mu-x}
        -
        \E[D_e(x\,\|\,X)]\,,
    \end{align}
    and the proof is completed by showing that $g=\E[\nabla e(X)]$ and
    $\tr(H\Sigma)=\E[\ip{\nabla e(X)}{X-\mu}]$.

    Since $e$ is Lipschitz, it is differentiable Lebesgue-almost everywhere,
    and its almost-everywhere partial derivatives are bounded by $\Lip(e)$
    and represent its weak derivatives.
    Applying this identity in every coordinate yields $g=\E[\nabla e(X)]$,
    which proves the first identity.

    We next identify the trace term. Since
    $H=\langle D^2e,\gamma_{\mu,\Sigma}\rangle$,
    distributional integration by parts gives
    \begin{align}
        \tr(H\Sigma)
        =
        \sum_{i,j=1}^d
        \Sigma_{ij}
        \left\langle
        \partial_{ij}e,
        \gamma_{\mu,\Sigma}
        \right\rangle
        =
        -\sum_{i,j=1}^d
        \Sigma_{ij}
        \left\langle
        \partial_i e,
        \partial_j\gamma_{\mu,\Sigma}
        \right\rangle
        =
        -\int_{\R^d}
        \left\langle
        \nabla e(x),
        \Sigma\nabla\gamma_{\mu,\Sigma}(x)
        \right\rangle
        \,dx.
    \end{align}
    The Gaussian score identity
    $\Sigma\nabla\gamma_{\mu,\Sigma}(x)
        =
        -(x-\mu)\gamma_{\mu,\Sigma}(x)$
    therefore yields
    \begin{align}
        \tr(H\Sigma)
         & =
        \int_{\R^d}
        \left\langle
        \nabla e(x),
        x-\mu
        \right\rangle
        \gamma_{\mu,\Sigma}(x)\,dx
        =
        \E\left[
            \left\langle
            \nabla e(X),
            X-\mu
            \right\rangle
            \right].
        \qedhere
    \end{align}
\end{proof}

\begin{corollary}[Regret switching]
    \label{cor:regret-switch}
    For
    $q(x)=\ip{g}{x-\mu}+\frac\lambda2\normm{x-\mu}{H}^2$,
    \begin{align}
        q(\mu) - q(x)
        =\bigl(\E[e(X)]-e(x)\bigr)
        -\tr(H\Sigma)
        -\frac\lambda2\normm{x-\mu}{H}^2
        +\E[D_e(x\,\|\,X)].
        \label{eq:q-regret-switch}
    \end{align}
\end{corollary}
Adding $\lambda\normm{x-\mu}{H}^2/2$ to both sides
of \eqref{eq:regret-switch} and the definition of $q$ proves \cref{cor:regret-switch}.
The left-hand side of \eqref{eq:q-regret-switch} can be seen as \textit{deterministic quadratic regret} $q(\mu) - q(x)$, while the right-hand side is the actual regret term $\E[e(X)]-e(x)$.
The following proposition replaces the lower bound supplied by optimistic
smoothing in the usual analysis.

\begin{proposition}
    \label{prop:bracket}
    Let
    $e:\R^d\to\R$ be convex and $L$-Lipschitz and $X\sim\N(\mu,\Sigma)$.
    Then for any $\beta\in(0,1/2]$ and $\lambda\in(0, 1/8]$
    such that $\lambda\left(d+\log\frac2\beta\right) \leq \frac{1}{64}$,
    for all $x\in\R^d$ satisfying
    \begin{align}
        \lambda\normm{\mu-x}{\Sigma^{-1}}
        \sqrt{\log\frac2\beta}                \leq \frac{1}{64},
        \label{eq:bracket-focus}
    \end{align}
    it holds that
    \begin{align}
        \frac\lambda2\normm{\mu-x}{H}^2
        -\E[D_e(x\,\|\,X)]
        \leq
        \frac1{16}\tr(H\Sigma)
        +\frac{\lambda L\beta}{\sqrt{\pi}}
        \norm{\mu-x}\normm{\mu-x}{\Sigma^{-1}}.
        \label{eq:bracket}
    \end{align}
\end{proposition}
The proof is given in \cref{app:bracket}.
% Its main idea is short.  After
% standardising the Gaussian, write the Bregman divergence as the curvature
% integrated along the segment from the comparator to the sample.
% On a high-probability set, the laws of the last $O(\lambda)$ portion of this segment
% dominate a constant fraction of the standard Gaussian law.  This pays for the
% comparator curvature.  One remaining part is bounded by $\tr(H\Sigma)$, while
% the Gaussian tail is controlled using Lipschitzness.

\subsection{Online Newton Step}
\begin{lemma}[Projected ONS]
    \label{lem:ons-step}
    On any round with $\Sigma_{t+1}\succ0$ in
    \cref{alg:direct-hermite}, every $x\in\K_\eps$ satisfies
    \begin{align}
        \frac12
        \left(\normm{\mu_{t+1}-x}{\Sigma_{t+1}^{-1}}^2
        -\normm{\mu_t-x}{\Sigma_t^{-1}}^2\right)
        \leq
        \eta\left(
        \widehat q_t(x)
        -
        \widehat q_t(\mu)
        \right)
        +\frac{\eta^2}{2}\normm{\widehat g_t}{\Sigma_{t+1}}^2
        \,.
        \label{eq:ons-step}
    \end{align}
\end{lemma}
\begin{proof}
    Since $\mu_{t+1}$ is the projection of $\mu_t-\eta\Sigma_{t+1}\widehat g_t$ onto
    $\K_\eps$ we have
    \begin{align}
        \frac12\normm{\mu_{t+1}-x}{\Sigma_{t+1}^{-1}}^2
         & \leq
        \frac12\normm{
            x - (
            \mu_t - \eta\Sigma_{t+1}\widehat g_t)
        }
        {\Sigma_{t+1}^{-1}}^2
        =
        \frac12\normm{\mu_t-x}{\Sigma_{t+1}^{-1}}^2
        +
        \frac{\eta^2}{2}\normm{g_t}{\Sigma_{t+1}^{-1}}^2
        + \eta\ip{\widehat g_t}{x-\mu_t}
        \\
         &
        =\frac12\normm{\mu_t-x}{\Sigma_t^{-1}}^2
        +\eta\ip{\widehat g_t}{x-\mu_t}
        +\frac{\lambda\eta}{2}(x-\mu_t)^\top\widehat H_t(x-\mu_t)
        +\frac{\eta^2}{2}\normm{\widehat g_t}{\Sigma_{t+1}}^2
        \\
         & =
        \frac12\normm{\mu_t-x}{\Sigma_t^{-1}}^2
        % +\eta\ip{\widehat g_t}{x-\mu_t}
        % +\frac{\lambda\eta}{2}(x-\mu_t)^\top\widehat H_t(x-\mu_t)
        +\frac{\eta^2}{2}\normm{\widehat g_t}{\Sigma_{t+1}}^2
        -
        \eta\left(
        \widehat q_t(\mu)
        -
        \widehat q_t(x)
        \right)\,,
    \end{align}
    which is \eqref{eq:ons-step}.
\end{proof}

\subsection{Stochastic control}
\label{sec:stochastic-control}
The random Hessian estimate $\hat{H}_t$ is indefinite, so $\Sigma_t$ need not decrease.
Introduce the predictable covariance
\begin{align}
    \overline\Sigma_1^{-1}=\Sigma_1^{-1} = \sigma^{-2}I_d,
    \qquad
    \overline\Sigma_{t+1}^{-1}
    =\overline\Sigma_t^{-1}+\lambda\eta H_t.
    \label{eq:predictable-covariance}
\end{align}
Since $H_t\succeq0$, the matrices $\overline\Sigma_t$ decrease.
The covariance stopping time is
\begin{align}
    \tau_{\mathrm{cov}}
    =
    \min\left(
    \left\{
    t \in [n]:
    \frac12 \overline\Sigma_{t+1}^{-1}
    \preceq
    \Sigma_t^{-1}+\lambda\eta\widehat H_t
    \preceq
    \frac32 \overline\Sigma_{t+1}^{-1}
    \text{ fails}
    \right\}
    \cup\{n\}
    \right).
    \label{eq:covariance-stopping-time}
\end{align}
Thus $\tau_{\mathrm{cov}}$ is the first covariance-comparison failure, or $n$
if the comparison never fails.  This is an $(\mathcal F_t)$-stopping time.
By minimality and the equality $\Sigma_1^{-1}=\overline\Sigma_1^{-1}$,
\begin{align}
    \frac12\overline\Sigma_t^{-1}
    \preceq\Sigma_t^{-1}
    \preceq\frac32\overline\Sigma_t^{-1}
    \qquad(1\leq t\leq\tau_{\mathrm{cov}}).
    \label{eq:covariance-comparison}
\end{align}
\begin{lemma}
    \label{lem:trace}
    % Suppose that $\sqrt{\frac32} \lambda\eta \Lip(e) \sigma \leq1$.
    Suppose that $\frac{c_\lambda c_\eta}{C_\eps C_0^6}\leq\frac{1}{8}$.
    Then it holds that
    \begin{align}
        \|\overline\Sigma_{t+1}^{-1}\|_{\op}^{\frac12}
        \leq
        \sigma^{-1}+t\,,
        % \lambda\eta\Lip(e)
        \qquad(1\leq t\leq \tau_{\mathrm{cov}}),
        \label{eq:predictable-spectrum}
    \end{align}
    and
    \begin{align}
        \eta\sum_{t=1}^{\tau_{\mathrm{cov}}}
        \tr(H_t\overline\Sigma_t)
        \leq\frac{4d\log(1 + n)}{\lambda}.
        \label{eq:trace-budget}
    \end{align}
\end{lemma}
The proof is in \cref{app:trace}.

\begin{lemma}
    \label{lem:feedback}
    Suppose that $C_0\geq 64$, $C_\eps> 2^{11}$, $c_\lambda c_\eta C_\eps> 2^{12}$, $\eps\leq 1/2$,
    $\ell\geq 1$, and $\frac{c_\lambda c_\eta}{C_\eps C_0^6}\leq\frac{1}{8}$.
    It holds that
    \begin{align}
        \sum_{t=1}^{\tau_{\mathrm{cov}}}\E_t[Y_t^2]\leq 6n\,,
        \label{eq:predictable-Y2}
    \end{align}
    and with probability at least $1-\delta/8$,
    \begin{align}
        \sum_{t=1}^{\tau_{\mathrm{cov}}}Y_t^2 \leq 8n.
        \label{eq:realised-Y2}
    \end{align}
\end{lemma}

The proof is in \cref{app:feedback}.

\subsection{Stability and regret}
\label{sec:stability}

Let
\begin{align}
    x_\star\in\argmin_{x\in\K_\eps}f(x).
    \label{eq:xstar}
\end{align}
This point exists because $\K_\eps$ is a compact subset of
$\operatorname{int}\K$, where the finite convex function $f$ is continuous.
Moreover, $x/\pi_\eps(x)\in\K_\eps$ for every $x\in\R^d$, so agreement on
$\K_\eps$ and \eqref{lem:extension} show that $x_\star$ is a global
minimizer of $e$.  Consequently,
\begin{align}
    \E_t[e(X_t)]-e(x_\star)\geq0.
    \label{eq:stationarity-sign}
\end{align}
This is the only point at which stationarity is essential.

Choose a sufficiently small universal constant $b>0$ and define the focus
radius
$B=\frac{b}{\lambda^2\ell}$
and the stopping times
\begin{align}
    \tau_{\mathrm{focus}}
    =
    \min\left(
    \left\{
    t\in[n]:
    \frac12
    \normm{\mu_{t+1}-x_\star}{\Sigma_{t+1}^{-1}}^2>B
    \right\}
    \cup\{n\}
    \right)
    \qquad\text{and}\qquad
    \tau_{\mathrm{safe}}
    =
    \tau_{\mathrm{cov}}\wedge\tau_{\mathrm{focus}}.
    \label{eq:safe-stopping-time}
\end{align}
Both are $(\F_t)$-stopping times.
Moreover, since $\mu_1=0$, $\Sigma_1^{-1}=\sigma^{-2}I_d$,
and $\norm{x_\star}\leq2d$,
\begin{align}
    \frac12\normm{\mu_1-x_\star}{\Sigma_1^{-1}}^2
    \leq
    4d^3M^2
    \leq
    \frac{b}{c_\lambda^2}d^3M^2\ell^3
    =B,
\end{align}
where the second inequality follows by taking $c_\lambda$ sufficiently small
relative to $b$.
Hence, by minimality of $\tau_{\mathrm{safe}}$,
\begin{align}
    \frac12\normm{\mu_t-x_\star}{\Sigma_t^{-1}}^2
    \leq B,
    \qquad
    \frac12\overline\Sigma_t^{-1}
    \preceq\Sigma_t^{-1}
    \preceq\frac32\overline\Sigma_t^{-1}
    \qquad(t\leq\tau_{\mathrm{safe}}).
    \label{eq:stability-stopped-invariants}
\end{align}
Choosing $b$ sufficiently small, the stopped focus bound also gives
\begin{align}
    \lambda\normm{\mu_t-x_\star}{\Sigma_t^{-1}}\sqrt\ell
    \leq\sqrt{2b}\leq c
    \qquad(t\leq\tau_{\mathrm{safe}}),
    \label{eq:concentration-focus}
\end{align}
where $c>0$ is a sufficiently small universal constant.

The next proposition contains the stochastic estimates used in both the
stability and regret arguments.  Its stopped formulation avoids any circular
use of covariance comparison or focus.
\begin{proposition}
    \label{prop:concentration}
    With probability at least $1-\delta/4$, simultaneously,
    \begin{align}
        \left|
        \sum_{t=1}^{\tau_{\mathrm{safe}}}
        \bigl(\widehat q_t(x_\star)-q_t(x_\star)\bigr)
        \right|
         & \leq
        \frac{C\sqrt{n\ell}}{\lambda},
        \label{eq:q-concentration}
        \\
        \sum_{t=1}^{\tau_{\mathrm{safe}}}
        \normm{\widehat g_t}{\Sigma_t}^2
         & \leq
        Cdn\ell,
        \label{eq:gradient-energy}
    \end{align}
    and
    \begin{align}
        -C\ell^{3/2}\sqrt{dn}\,
        \overline\Sigma_{\tau_{\mathrm{safe}}+1}^{-1}
        \preceq
        \sum_{t=1}^{\tau_{\mathrm{safe}}}
        (\widehat H_t-H_t)
        \preceq
        C\ell^{3/2}\sqrt{dn}\,
        \overline\Sigma_{\tau_{\mathrm{safe}}+1}^{-1}.
        \label{eq:H-concentration}
    \end{align}
\end{proposition}

The proof is in \cref{app:concentration}.

\begin{proposition}
    \label{prop:stability}
    With probability at least $1-\delta/2$, $\tau_{\mathrm{safe}}=n$.
\end{proposition}

The proof is in \cref{app:stability}.

\begin{proof}[Proof of \cref{thm:main}]
    If $\eps_0>1/2$, then
    \begin{align}
        \operatorname{Reg}_n
        \leq n
        =\sqrt n\sqrt n
        \leq2C_\eps d^2M\ell^4\sqrt n,
    \end{align}
    so \eqref{eq:main-regret} holds.  Suppose henceforth that
    $\eps=\eps_0\leq1/2$.

    Work on the event of \cref{prop:stability}.  The focus bound gives
    $\lambda\normm{\mu_t-x_\star}{\Sigma_t^{-1}}\sqrt\ell\leq\sqrt{2b}$.
    Since $d^{3/2}M\geq d$ and $\log(2e^\ell)\leq2\ell$,
    \begin{align}
        \lambda\left(d+\log(2e^\ell)\right)
         & \leq Cc_\lambda\leq c_0,
         &
        \lambda\normm{\mu_t-x_\star}{\Sigma_t^{-1}}
        \sqrt{\log(2e^\ell)}
         & \leq C\sqrt b\leq c_0.
    \end{align}
    We also have $\lambda\leq1/8$ and $e^{-\ell}<1/2$, so
    \cref{prop:bracket} applies.  Moreover,
    $\mu_t,x_\star\in\K_\eps\subseteq B_{2d}$ gives
    \begin{align}
        \sum_{t=1}^n
        \lambda\Lip(e)e^{-\ell}
        \norm{\mu_t-x_\star}
        \normm{\mu_t-x_\star}{\Sigma_t^{-1}}
        \leq
        4dn\lambda\Lip(e)e^{-\ell}\sqrt{2B}
        =
        \frac{Cdn\Lip(e)e^{-\ell}}{\sqrt\ell}
        \leq
        \frac{Cd n^{3/2}e^{-\ell}}{\sqrt\ell}
        \leq C.
    \end{align}
    Starting from the cumulative expected extension regret, applying
    \cref{cor:regret-switch,prop:bracket}, inserting
    $q_t(x_\star)=\widehat q_t(x_\star)
        -(\widehat q_t(x_\star)-q_t(x_\star))$, and summing
    \eqref{eq:ons-step} with $\widehat q_t(\mu_t)=0$ gives
    \begin{align}
         & \sum_{t=1}^n
        \left(\E_t[e(X_t)]-e(x_\star)\right)
        \nonumber                        \\
         & \quad\leq
        -\sum_{t=1}^n q_t(x_\star)
        +C\sum_{t=1}^n\tr(H_t\Sigma_t)
        +C
        \nonumber                        \\
         & \quad=
        -\sum_{t=1}^n\widehat q_t(x_\star)
        +\sum_{t=1}^n
        \bigl(\widehat q_t(x_\star)-q_t(x_\star)\bigr)
        +C\sum_{t=1}^n\tr(H_t\Sigma_t)
        +C
        \nonumber                        \\
         & \quad\leq
        \frac{\normm{\mu_1-x_\star}{\Sigma_1^{-1}}^2}{2\eta}
        +\frac\eta2\sum_{t=1}^n
        \normm{\widehat g_t}{\Sigma_{t+1}}^2
        +\left|\sum_{t=1}^n
        \bigl(\widehat q_t(x_\star)-q_t(x_\star)\bigr)\right|
        \nonumber                        \\
         & \qquad
        +C\sum_{t=1}^n\tr(H_t\Sigma_t)+C
        \label{eq:regret-decomposition}
        \\
         & \quad\leq
        \frac{\normm{\mu_1-x_\star}{\Sigma_1^{-1}}^2}{2\eta}
        +C\eta dn\ell
        +\frac{C\sqrt{n\ell}}{\lambda}
        +\frac{Cd\ell}{\lambda\eta}
        +C
        \nonumber                        \\
         & \quad\leq
        Cd^2M\sqrt n
        +Cd^2M\ell\sqrt n
        +Cd^{3/2}M\ell^{5/2}\sqrt n
        +Cd^{3/2}\ell^3\sqrt n
        +C
        \nonumber                        \\
         & \quad\leq Cd^2M\ell^3\sqrt n.
        \label{eq:expected-extension-regret}
    \end{align}
    Here the fourth line uses \eqref{eq:next-gradient-energy},
    \eqref{eq:q-concentration}, and \eqref{eq:stability-stopped-trace}; the fifth
    follows by substituting \eqref{eq:parameters-one}--\eqref{eq:parameters-two},
    using $\mu_1=0$ and
    $\norm{x_\star}\leq2d$; the last uses $M\geq d^{-1/2}$ and $\ell\geq1$.

    We now return to the original loss.  For every $x\in\K$, the point
    $(1-\eps)x$ lies in $\K_\eps$, so convexity and $f\in[0,1]$ give
    \begin{align}
        f(x_\star)
        \leq f((1-\eps)x)
        \leq(1-\eps)f(x)+\eps f(0)
        \leq f(x)+\eps.
    \end{align}
    Taking the infimum over $x$ yields
    \begin{align}
        f(x_\star)-\inf_{x\in\K}f(x)\leq\eps.
        \label{eq:shrink-comparator}
    \end{align}
    Moreover, $A_t\in\K_\eps$, and therefore
    \begin{align}
        f(A_t)=e(A_t)\leq e(X_t)
    \end{align}
    by \eqref{lem:extension}.  Since $f(A_t)\in[0,1]$, Azuma's
    inequality gives, with probability at least $1-\delta/2$,
    \begin{align}
        \sum_{t=1}^n
        \left(f(A_t)-\E_t[f(A_t)]\right)
        \leq C\sqrt{n\ell}
        \label{eq:action-martingale}
    \end{align}
    Intersecting this event with the event of \cref{prop:stability} and
    starting from the realized regret, we obtain
    \begin{align}
        \operatorname{Reg}_n
         & =
        \sum_{t=1}^n
        \left(f(A_t)-\E_t[f(A_t)]\right)
        \nonumber                   \\
         & \qquad
        +\sum_{t=1}^n
        \left(\E_t[f(A_t)]-f(x_\star)\right)
        +n\left(f(x_\star)-\inf_{x\in\K}f(x)\right)
        \nonumber                   \\
         & \leq
        C\sqrt{n\ell}+n\eps
        +\sum_{t=1}^n
        \left(\E_t[e(X_t)]-e(x_\star)\right)
        \\
         & \leq Cd^2M\ell^4\sqrt n,
    \end{align}
    where the last inequality uses \eqref{eq:expected-extension-regret} and
    $n\eps=C_\eps d^2M\ell^4\sqrt n$.  The two events fail with total
    probability at most $\delta$, completing the proof.
\end{proof}

\newpage
\begin{thebibliography}{9}
    \raggedright

    \bibitem[Freedman(1975)]{freedman1975tail}
    David A. Freedman.
    \newblock On tail probabilities for martingales.
    \newblock \emph{The Annals of Probability}, 3(1):100--118, 1975.

    \bibitem[Lattimore(2026)]{lattimore2026bandit}
    Tor Lattimore.
    \newblock \emph{Bandit Convex Optimisation}.
    \newblock Cambridge University Press, 2026.
    \newblock \href{https://doi.org/10.1017/9781009607551}{doi:10.1017/9781009607551}.

    \bibitem[Lattimore and Gy\"orgy(2023)]{lattimore2023second}
    Tor Lattimore and Andr\'as Gy\"orgy.
    \newblock A second-order method for stochastic bandit convex optimisation.
    \newblock In \emph{Proceedings of the 36th Conference on Learning Theory},
    volume 195 of \emph{Proceedings of Machine Learning Research}, pages
    2067--2094, 2023.

    \bibitem[Vershynin(2018)]{vershynin2018high}
    Roman Vershynin.
    \newblock \emph{High-Dimensional Probability: An Introduction with
        Applications in Data Science}.
    \newblock Cambridge University Press, 2018.


    \bibitem[Evans and Gariepy(2015)]{evansgariepy2015measure}
    Lawrence C. Evans and Ronald F. Gariepy.
    \newblock \emph{Measure Theory and Fine Properties of Functions}.
    \newblock Revised edition, CRC Press, 2015.

\end{thebibliography}

\newpage
\appendix
\section{Hessian Measures}
First we record the standard fact that the distributional Hessian of a convex function is a positive-semidefinite matrix-valued Radon measure.
\begin{lemma}[Hessian measure of a convex function]
    \label{lem:convex-hessian-measure}
    Let $r:\R^d\to\R$ be convex and Lipschitz.
    Then the distributional Hessian $D^2r$ is a positive-semidefinite
    matrix-valued Radon measure.
\end{lemma}
\begin{proof}
    This is standard; see \citet[Section~6.3]{evansgariepy2015measure}.
\end{proof}
Next lemma gives the Gaussian averaging identity for the Hessian measure, which is the key to the direct Hermite coefficients.
\begin{lemma}[Gaussian averaging of the Hessian measure]
    \label{lem:gaussian-hessian-average}
    Let $r:\R^d\to\R$ be convex and Lipschitz, let
    $X=\mu+\Sigma^{1/2}Z$, where $Z\sim\N(0,I_d)$ and $\Sigma\succ0$, and
    let $\gamma_{\mu,\Sigma}$ be the density of $X$.  Then
    \begin{align}
        \langle D^2r,\gamma_{\mu,\Sigma}\rangle
        =
        \int_{\R^d}
        \gamma_{\mu,\Sigma}(x)\,D^2r(dx)
        =
        \E\left[
            r(X)\Sigma^{-1/2}(ZZ^\top-I_d)\Sigma^{-1/2}
            \right].
        \label{eq:gaussian-hessian-average}
    \end{align}
\end{lemma}
\begin{proof}
    By \cref{lem:convex-hessian-measure}
    $D^2r$ is a positive-semidefinite matrix-valued Radon measure.
    Since $r$ is Lipschitz, $D^2r \in \mathcal S'(\R^d)$ is a tempered distribution,
    and the Gaussian density $\gamma_{\mu,\Sigma}\in\mathcal S(\R^d)$ is a Schwartz function,
    which means that $\langle D^2r,\gamma_{\mu,\Sigma}\rangle < \infty$ and
    \begin{align}
        \langle D^2r,\gamma_{\mu,\Sigma}\rangle
        =
        \langle r,D^2\gamma_{\mu,\Sigma}\rangle
        =
        \E\left[
            r(X)\Sigma^{-1/2}(ZZ^\top-I_d)\Sigma^{-1/2}
            \right],
    \end{align}
    where the first equality is integration by parts applied twice.
\end{proof}

\section{Proof \cref{prop:bracket}}
\label{app:bracket}
Let $M$ be a positive-semidefinite matrix-valued Radon measure and let
$\nu=\tr M$ be its trace measure.
Positivity implies that every entry of $M$ is absolutely continuous with
respect to $\nu$.\alireza{Needs proof or citation.}
Hence, by the Radon--Nikodym theorem, there is a
measurable matrix-valued function $Q$ such that\alireza{Needs proof.}
\begin{align}
    M(dz)=Q(z)\nu(dz),
    \qquad
    Q(z)\succeq0,
    \qquad
    \tr Q(z)=1
    \quad\text{for $\nu$-almost every }z.
    \label{eq:matrix-measure-polar}
\end{align}
Therefore, for a nonnegative measurable function $\rho$ and a measurable vector field
$b$, we use the notation
\begin{align}
    \int_B \rho(z)b(z)^\top M(dz)b(z)
    \triangleq
    \int_B
    \rho(z)b(z)^\top Q(z)b(z)\,\nu(dz).
    \label{eq:matrix-measure-quadratic-form}
\end{align}
In particular, this integral is nonnegative, possibly with value $+\infty$.

\begin{lemma}[Averaged Bregman curvature]
    \label{lem:averaged-bregman-curvature}
    Let $h:\R^d\to\R$ be convex and Lipschitz, let $\bar{x}\in\R^d$, and let
    $Z\sim\N(0,I_d)$.  For $s\in(0,1]$, let $p_s$ denote the density of $Z_s=(1-s)\bar{x}+sZ$.
    Then
    \begin{align}
        \E[D_h(\bar{x}\,\|\,Z)]
        =
        \int_0^1\frac1s
        \int_{\R^d}
        p_s(z)(z-\bar{x})^\top D^2h(dz)(z-\bar{x})\,ds.
        \label{eq:averaged-bregman-curvature}
    \end{align}
\end{lemma}
\begin{proof}
    TODO. It is just FTC.
\end{proof}


\begin{proof}[Proof of \cref{prop:bracket}]
    First, we switch to the standardized case by defining
    \begin{align}
        \bar{x}=\Sigma^{-\frac12}(x-\mu),
        \qquad
        h(z)=e(\mu+\Sigma^{1/2}z),
        \quad\text{and}
        \quad
        Z=\Sigma^{-\frac12}(X - \mu) \sim \N(0,I_d)\,.
    \end{align}
    Finally, since $\Sigma^{\frac12}\bar{x}=x-\mu$, the chain rule at
    differentiability points of $e$ gives
    $D_h(\bar{x}\,\|\,Z)=D_e(x\,\|\,X)$
    almost surely.

    The plan is to define a set $\mathcal A\subset \R^d$ and control the integrals over $\mathcal A$ and $\mathcal A^c$ separately.
    To this end, write
    \begin{align}
        \normm{\mu-x}{H}^2
        =
        \bar{x}^\top \Sigma^{\frac12}H\Sigma^{\frac12}\bar{x}
         & =
        \int_{\R^d}\gamma(z)\bar{x}^\top D^2h(dz)\bar{x}
        =
        \int_{\mathcal A}\gamma(z)\bar{x}^\top D^2h(dz)\bar{x}
        +
        \int_{\mathcal A^c}\gamma(z)\bar{x}^\top D^2h(dz)\bar{x}
        \label{eq:standardized-comparator-curvature}
    \end{align}
    In step 1, we upper bound the integral over $\mathcal A$ by the sum of expected Bregman divergence $\E[D_e(x\,\|\,X)]$ and a trace term $\tr(H\Sigma)$.
    In step 2, we upper bound the integral over $\mathcal A^c$ by a tail term that is small under the assumptions of the proposition.

    \proofstepinline{1}
    Let $\gamma=\gamma_{0,I_d}$ and $\gamma_V=\gamma_{0,V}$ for $V\succ0$.
    For $s\in(0,1]$ by \cref{lem:averaged-bregman-curvature},
    we have
    \begin{align}
        \E[D_e(x\,\|\,X)]
        =
        \E[D_h(\bar{x}\,\|\,Z)]
        =
        \int_0^1\frac1s
        \int_{\R^d}
        p_s(z)(z-\bar{x})^\top D^2h(dz)(z-\bar{x})\,ds\,,
        \label{eq:bregman-density}
    \end{align}
    where $p_s$ denotes the density of $(1-s)\bar{x}+sZ$.
    Define the set\footnote{
        If $\bar{x}=0$, then the left-hand side of \eqref{eq:bracket} is nonpositive, so the inequality holds trivially.
    }
    \begin{align}
        \mathcal A=\left\{z:
        \norm{z}^2\leq2d\log2+4\log\frac2\beta
        \quad
        \text{and}
        \quad
        \left|\ip{\frac{\bar{x}}{\norm{\bar{x}}}}{z}\right|^2
        \leq4\log\frac{2\sqrt2}{\beta}
        \right\}.
        \label{eq:A}
    \end{align}

    \proofstepinline{1.1}
    Next we show that $\frac12\gamma(z)\leq p_s(z)$ for $s\in[1-2\lambda,1]$ and $z\in\mathcal A$.
    Since $sZ+(1-s)\bar{x}\sim\N((1-s)\bar{x},s^2I_d)$,
    we have
    \begin{align}
        \log\frac{p_s(z)}{\gamma(z)}
         & =-d\log s
        -\frac{1-s^2}{2s^2}\norm{z}^2
        +\frac{1-s}{s^2}\ip{\bar{x}}{z}
        -\frac{(1-s)^2}{2s^2}\norm{\bar{x}}^2\,.
        \label{eq:density-ratio}
    \end{align}
    For $s\in[1-2\lambda,1]$, the assumption $\lambda\leq1/8$ implies $s\geq\frac34$, and $0\leq1-s\leq2\lambda$.
    Consequently,
    \begin{align}
        \frac{1-s^2}{2s^2}
        \leq
        \frac{1-s}{s^2}
        \leq
        \frac{2\lambda}{(3/4)^2}
        =
        \frac{32}{9}\lambda
        \leq4\lambda,
        \quad\text{and}\quad
        \frac{(1-s)^2}{2s^2}
        \leq
        \frac{(2\lambda)^2}{2(3/4)^2}
        =
        \frac{32}{9}\lambda^2
        \leq4\lambda^2.
    \end{align}
    Moreover, $-d\log s\geq0$. Therefore,
    \eqref{eq:density-ratio} yields
    \begin{align}
        \log\frac{p_s(z)}{\gamma(z)}
        \geq
        -4\lambda\norm{z}^2
        -4\lambda\left|\ip{\bar{x}}{z}\right|
        -4\lambda^2\norm{\bar{x}}^2.
        \label{eq:density-ratio-lower}
    \end{align}
    Now suppose for $z\in\mathcal A$, then
    together with the assumption $\lambda\left(d+\log\frac2\beta\right)\leq \frac{1}{64}$ we have
    \begin{align}
        \lambda\norm{z}^2
         & \leq
        \lambda\left(
        2d\log2+4\log\left(\frac2\beta\right)
        \right)
        \leq
        4\lambda\left(d+\log\left(\frac2\beta\right)\right)
        \leq
        \frac{1}{16}\,,
        \label{ineq:lambda-z-squared-bound}
    \end{align}
    and, again, from $z\in\mathcal A$ we get
    \begin{align}
        \lambda\left|\ip{\bar{x}}{z}\right|
        \leq
        2\lambda\norm{\bar{x}}
        \sqrt{\log\frac{2\sqrt2}{\beta}}
        \leq
        2\lambda\norm{\bar{x}}
        \sqrt{
            \log\left(
            \frac{2}{\beta}
            \right)
            +\frac12\log2
        }
        \leq
        \sqrt5\,
        \lambda\norm{\bar{x}}\sqrt{\log\left(\frac{2}{\beta}\right)}
        \leq
        \frac{\sqrt5}{64}\,,
    \end{align}
    where in the penultimate inequality we used $
        \frac12\log2
        \leq \frac14\log\left(\frac{2}{\beta}\right)$ since $\beta<1/2$,
    and in the last inequality we used \eqref{eq:bracket-focus}.
    The last inequality also gives
    $\lambda^2\norm{\bar{x}}^2
        \leq
        \frac{1}{2^{12} \log\left(\frac{2}{\beta}\right)}
        \leq
        \frac{1}{2^{12} \log4}$.
    Substituting these three bounds into
    \eqref{eq:density-ratio-lower} gives
    \begin{align}
        \log\frac{p_s(z)}{\gamma(z)}
        \geq
        -4\left(\frac{1}{16}+\frac{\sqrt5}{64}+\frac{1}{2^{12} \log4}\right)
        \geq
        \log\frac12\,.
    \end{align}
    Therefore, we have established
    $p_s(z)
        \geq
        \frac12\gamma(z)$
    for all
    $s\in[1-2\lambda,1]$
    and
    $z\in\mathcal A$.

    \proofstepinline{1.2}
    Since the measure integrand in \eqref{eq:bregman-density} is
    nonnegative, we may restrict both integrals to
    $s\in[1-2\lambda,1]$ and $z\in\mathcal A$,
    which together with $-\log(1-2\lambda)\geq2\lambda$ gives
    \begin{align}
        \E[D_h(\bar{x}\,\|\,Z)]
         & \geq
        \int_{1-2\lambda}^1\frac1s
        \int_{\mathcal A}
        p_s(z)(z-\bar{x})^\top D^2h(dz)(z-\bar{x})\,ds
        \\
         & \geq
        \frac12
        \int_{1-2\lambda}^1\frac{ds}{s}
        \int_{\mathcal A}
        \gamma(z)(z-\bar{x})^\top D^2h(dz)(z-\bar{x})
        \\
         & =
        \frac{-\log(1-2\lambda)}{2}
        \int_{\mathcal A}
        \gamma(z)(z-\bar{x})^\top D^2h(dz)(z-\bar{x})
        \\
         & \geq
        \lambda
        \int_{\mathcal A}
        \gamma(z)(z-\bar{x})^\top D^2h(dz)(z-\bar{x}).
        \label{eq:bregman-lower}
    \end{align}
    Apply \eqref{eq:matrix-measure-polar} to $D^2h$ and write
    $\nu=\tr(D^2h)$, and $D^2h(dz)=Q(z)\nu(dz)$,
    where $Q(z)\succeq0$ and $\tr Q(z)=1$ for $\nu$-almost every $z$.
    For every positive-semidefinite matrix $G$, by the triangle inequality we have
    \begin{align}
        \frac12\bar{x}^\top G\bar{x}
        \leq
        (z-\bar{x})^\top G(z-\bar{x})+z^\top Gz.
    \end{align}
    Applying this inequality to $G=Q(z)$, multiplying by
    $\lambda\gamma(z)$, and integrating over $\mathcal A$ gives
    \begin{align}
        \frac\lambda2
        \int_{\mathcal A}
        \gamma(z)\bar{x}^\top D^2h(dz)\bar{x}
        \leq{} &
        \lambda
        \int_{\mathcal A}
        \gamma(z)(z-\bar{x})^\top D^2h(dz)(z-\bar{x})
        +
        \lambda
        \int_{\mathcal A}
        \gamma(z)z^\top D^2h(dz)z.
    \end{align}
    Moreover, for $\nu$-almost every $z$,
    $z^\top Q(z)z
        \leq
        \norm{z}^2\tr Q(z)
        =
        \norm{z}^2.$
    On $\mathcal A$, \eqref{ineq:lambda-z-squared-bound}
    and letting $c_0=\frac{1}{64}$ gives $\lambda\norm{z}^2\leq1/16$.
    Hence,
    \begin{align}
        \lambda
        \int_{\mathcal A}
        \gamma(z)z^\top D^2h(dz)z
        \leq
        \frac1{16}
        \int_{\mathcal A}\gamma(z)\,\nu(dz)
        \leq
        \frac1{16}
        \int_{\R^d}\gamma(z)\tr(D^2h)(dz)
        =
        \frac1{16}\tr(H\Sigma),
        \label{eq:trace-upper-bound}
    \end{align}
    where the last equality is followed from
    \begin{align}
        \int_{\R^d}\gamma(z)\,\tr(D^2h)(dz)
        %  & =
        % \tr(\langle D^2h, \gamma\rangle)
        =
        \tr(\langle h, D^2\gamma\rangle)
        =
        \tr\left(\E\left[h(Z)(ZZ^\top-I_d)\right]\right)
        =
        \tr(\Sigma^{\frac12}H\Sigma^{\frac12})
        =
        \tr(H\Sigma).
        \label{eq:standardized-trace-curvature}
    \end{align}
    Finally, combining \cref{eq:bregman-lower,eq:bregman-lower,eq:trace-upper-bound} we obtain
    \begin{align}
        \frac\lambda2
        \int_{\mathcal A}
        \gamma(z)\bar{x}^\top D^2h(dz)\bar{x}
        \leq
        \E[D_h(\bar{x}\,\|\,Z)]
        +
        \frac1{16}\tr(H\Sigma)\,,
        \label{eq:bracket-step1}
    \end{align}
    which concludes step 1.

    \proofstepinline{2}
    It remains to bound the curvature measure on $\mathcal A^c$.  Set
    \begin{align}
        u=\frac{\bar{x}}{\norm{\bar{x}}},
        \qquad
        D^2_{\bar{x}}\triangleq\bar{x}^\top D^2h\,\bar{x}.
    \end{align}
    Since $D^2h$ is positive semidefinite, $D^2_{\bar{x}}$ is a nonnegative Radon
    measure.\alireza{we could've just used $D_{\bar{x}}^2$ from the beginning.}
    Direct comparison of Gaussian densities gives
    \begin{align}
        \gamma(z)\1\left\{\norm z^2>
        2d\log2+4\log\frac2\beta\right\}
        \leq
        \frac\beta2\gamma_{2I_d}(z)\,\,\,\text{and}\,\,\,
        \gamma(z)\1\left\{|\ip{u}{z}|^2>
        4\log\frac{2\sqrt2}{\beta}\right\}
        \leq
        \frac\beta2\gamma_{I_d+uu^\top}(z).
        \label{eq:directional-tail}
    \end{align}
    Since $\mathcal A^c$ is contained in the union of the two events above
    and $D^2_{\bar{x}}$ is nonnegative, integrating these inequalities gives
    \begin{align}
        \int_{\mathcal A^c}\gamma(z)\,D^2_{\bar{x}}(dz)
        \leq
        \frac\beta2
        \left(
        \int_{\R^d}\gamma_{2I_d}(z)\,D^2_{\bar{x}}(dz)
        +
        \int_{\R^d}\gamma_{I_d+uu^\top}(z)\,D^2_{\bar{x}}(dz)
        \right).
        \label{eq:tail-curvature-domination}
    \end{align}
    We next bound each integral on the right-hand side of
    \eqref{eq:tail-curvature-domination}.  Fix $V\in\{2I_d,I_d+uu^\top\}$.
    Since $D^2_{\bar{x}}=\bar{x}^\top D^2h\,\bar{x}$, distributional integration by parts against the Schwartz function
    $\gamma_V$ gives
    \begin{align}
        \int_{\R^d}\gamma_V(z)\,D^2_{\bar{x}}(dz)
         & =
        -\int_{\R^d}
        \ip{\nabla h(z)}{\bar{x}}
        \ip{\nabla\gamma_V(z)}{\bar{x}}\,dz
        =
        \int_{\R^d}
        \ip{\nabla h(z)}{\bar{x}}
        \ip{V^{-1}\bar{x}}{z}
        \gamma_V(z)\,dz,
        \label{eq:wider-gaussian-ibp}
    \end{align}
    where we used
    $\nabla\gamma_V(z)=-V^{-1}z\,\gamma_V(z)$.
    At every differentiability point of $e$, the chain rule gives
    \begin{align}
        \ip{\nabla h(z)}{\bar{x}}
         & =
        \ip{\nabla e(\mu+\Sigma^{1/2}z)}
        {\Sigma^{1/2}\bar{x}}
        =
        \ip{\nabla e(\mu+\Sigma^{1/2}z)}
        {x-\mu}.
    \end{align}
    Since $e$ is $L$-Lipschitz, it follows that
    \begin{align}
        \left|\ip{\nabla h(z)}{\bar{x}}\right|
        \leq
        L\norm{x-\mu}
        \qquad\text{for almost every }z.
    \end{align}
    Therefore,
    \begin{align}
        \int_{\R^d}\gamma_V(z)\,D^2_{\bar{x}}(dz)
         & \leq
        L\norm{x-\mu}
        \int_{\R^d}
        \left|\ip{V^{-1}\bar{x}}{z}\right|
        \gamma_V(z)\,dz
        =
        L\norm{x-\mu}
        \sqrt{\frac{2}{\pi}}
        \sqrt{\bar{x}^\top V^{-1}\bar{x}}.
        \label{eq:wider-gaussian-bound}
    \end{align}
    For $V=2I_d$, we have
    $\bar{x}^\top V^{-1}\bar{x}=\frac12\norm{\bar{x}}^2$.
    For $V=I_d+uu^\top$, since $\norm u=1$,
    we have
    $V^{-1}
        =
        I_d-\frac12uu^\top$.
    Moreover, $\bar{x}=\norm{\bar{x}}\,u$, and therefore
    $\bar{x}^\top V^{-1}\bar{x}=
        \bar{x}^\top
        \left(I_d-\frac12uu^\top\right)\bar{x}
        =\frac12\norm{\bar{x}}^2$.
    Thus, for either
    $V\in\{2I_d,I_d+uu^\top\}$,
    we have
    $\sqrt{\bar{x}^\top V^{-1}\bar{x}} = \frac{\norm{\bar{x}}}{\sqrt2}$.
    Substituting this into \eqref{eq:wider-gaussian-bound} gives
    \begin{align}
        \int_{\R^d}\gamma_V(z)\,D^2_{\bar{x}}(dz)
        \leq
        \frac{L}{\sqrt\pi}
        \norm{x-\mu}\norm{\bar{x}}.
        \label{eq:wider-gaussian}
    \end{align}
    Combining \eqref{eq:tail-curvature-domination} and
    \eqref{eq:wider-gaussian} now yields
    \begin{align}
        \frac{\lambda}{2}\int_{\mathcal A^c}\gamma(z)\,D^2_{\bar{x}}(dz)
        \leq
        \frac{\beta L \lambda}{2\sqrt\pi}
        \norm{x-\mu}\norm{\bar{x}}.
        \label{eq:bracket-step2}
    \end{align}
    Finally, the definition of $\bar{x}$ gives
    $\norm{\bar{x}}^2=\normm{\mu-x}{\Sigma^{-1}}^2$,
    completing the second step.
    Plugging \eqref{eq:bracket-step1} and \eqref{eq:bracket-step2}
    into \eqref{eq:standardized-comparator-curvature}
    completes the proof.
    %     \frac\lambda2\normm{\mu-x}{H}^2
    %     -
    %     \E[D_e(x\,\|\,X)]
    %     \leq
    %     \frac1{16}\tr(H\Sigma)
    %     +
    %     \frac{\lambda L\beta}{2\sqrt\pi}
    %     \norm{\mu-x}
    %     \normm{\mu-x}{\Sigma^{-1}}.
    % \end{align}
    % Since $1/(2\sqrt\pi)$ is a universal constant, this proves
    % \eqref{eq:bracket}.
\end{proof}

\section{Proof of \cref{lem:trace}}
\label{app:trace}

\begin{proof}[Proof of \cref{lem:trace}]
    Let $L=\Lip(e)$.
    To prove \eqref{eq:predictable-spectrum}, write
    \begin{align}
        \norm{\overline\Sigma_{t+1}^{-1}}_{\op}
        =
        \norm{\overline\Sigma_{t}^{-1} + \lambda\eta H_t}_{\op}
        \leq
        \norm{\overline\Sigma_{t}^{-1}}_{\op} + \lambda\eta\norm{H_t}_{\op}\,.
    \end{align}
    Next, distributional integration by parts for every $v\in\mathbb{S}_1^d$ gives
    \begin{align}
        \norm{H_t}_{\op}
         & =
        \sup_{\|v\|=1} v^\top H_t v
        =
        \sup_{\|v\|=1}
        \int_{\R^d}e(x) D^2_v\gamma_{\mu_t,\Sigma_t}(x)\,dx
        =
        \sup_{\|v\|=1}
        -\int_{\R^d}\partial_ve(x)\,
        \partial_v\gamma_{\mu_t,\Sigma_t}(x)\,dx
        \\
         & =
        \sup_{\|v\|=1}
        \E_t\left[
            \ip{\nabla e(X_t)}{v}
            \ip{\Sigma_t^{-1}(X_t-\mu_t)}{v}
            \right]
        \leq
        \sup_{\|v\|=1}
        L \|v\|_{\Sigma_t^{-1}}
        =
        L \norm{\Sigma_t^{-1}}_{\op}^{1/2},
        \label{eq:H-quadratic-bound}
    \end{align}
    where the inequality is by Cauchy--Schwarz and the Lipschitz property of $e$.
    Plugging this bound into the previous inequality gives
    \begin{align}
        \norm{\overline\Sigma_{t+1}^{-1}}_{\op}
         & \leq
        \norm{\overline\Sigma_{t}^{-1}}_{\op} + \lambda\eta L \norm{\Sigma_t^{-1}}_{\op}^{1/2}
        \leq
        \norm{\overline\Sigma_{t}^{-1}}_{\op} + \sqrt{\frac{3}{2}}\lambda\eta L\norm{\overline\Sigma_{t}^{-1}}_{\op}^{1/2}
        \leq
        \left(\norm{\overline\Sigma_{t}^{-1}}_{\op}^{1/2} + \lambda\eta L\right)^2\,,
        \label{eq:predictable-spectrum-bound}
    \end{align}
    where the second inequality uses \eqref{eq:covariance-comparison}.
    Now, by taking square roots and iterating, we get
    \begin{align}
        \norm{\overline\Sigma_{t+1}^{-1}}_{\op}^{1/2}
        \leq
        \norm{\overline\Sigma_1^{-1}}_{\op}^{1/2} + t\lambda\eta L
        = \sigma^{-1}+t\lambda\eta L
        \leq \sigma^{-1} + t \frac{c_\lambda}{d^{3/2}M\ell^2} \cdot \frac{c_\eta dM}{\sqrt n} \cdot \frac{8}{\eps}
        \leq
        \sigma^{-1} + 8t\cdot \frac{c_\lambda c_\eta}{C_\eps C_0^6}
        \leq
        \sigma^{-1} + t\,,
    \end{align}
    which proves \eqref{eq:predictable-spectrum}.


    We next prove \eqref{eq:trace-budget}.  Starting from its left-hand side
    and using cyclicity of the trace gives
    \begin{align}
        \eta\sum_{t=1}^{\tau_{\mathrm{cov}}}
        \tr(H_t\overline\Sigma_t)
        =
        \frac1\lambda\sum_{t=1}^{\tau_{\mathrm{cov}}}
        \tr\left(
        \lambda\eta\overline\Sigma_t^{1/2}H_t
        \overline\Sigma_t^{1/2}
        \right).
        \label{eq:normalised-trace-sum}
    \end{align}
    For every $t\leq\tau_{\mathrm{cov}}$ and every unit vector $u$, the same integration by
    parts estimate used above, followed by
    \eqref{eq:covariance-comparison}, gives
    \begin{align}
        u^\top\overline\Sigma_t^{1/2}H_t
        \overline\Sigma_t^{1/2}u
         & \leq
        L\norm{\overline\Sigma_t^{1/2}u}
        \sqrt{u^\top\overline\Sigma_t^{1/2}
            \Sigma_t^{-1}\overline\Sigma_t^{1/2}u}
        \leq
        \sqrt{\frac32}\,L
        \norm{\overline\Sigma_t^{1/2}u}
        \leq
        \sqrt{\frac32}\,L\sigma.
        \label{eq:normalised-H}
    \end{align}
    The last inequality uses
    $\overline\Sigma_t\preceq\overline\Sigma_1=\sigma^2I_d$.
    Therefore, by the lemma assumption
    \begin{align}
        \norm{\eta\lambda\overline\Sigma_t^{1/2}H_t\overline\Sigma_t^{1/2}}_{\op}
        \leq
        \sqrt{\frac32}\lambda\eta L\sigma
        \leq
        \sqrt{\frac32}\cdot\frac{c_\lambda}{d^{3/2}M\ell^2}\cdot\frac{c_\eta dM}{\sqrt n}\cdot\frac{8}{\eps}
        \cdot\frac1{M\sqrt{2d}}
        \leq
        8\cdot \frac{c_\lambda c_\eta}{C_\eps C_0^6}
        \leq
        1\,.
    \end{align}
    Returning to
    \eqref{eq:normalised-trace-sum} and applying
    $s\leq2\log(1+s)$ to eigenvalues of $\lambda\eta\overline\Sigma_t^{1/2}H_t\overline\Sigma_t^{1/2}$
    which are in $[0,1]$, we obtain
    \begin{align}
        \eta\sum_{t=1}^{\tau_{\mathrm{cov}}}
        \tr(H_t\overline\Sigma_t)
         & \leq
        \frac2\lambda\sum_{t=1}^{\tau_{\mathrm{cov}}}
        \log\det\left(
        I_d+\lambda\eta\overline\Sigma_t^{1/2}H_t
        \overline\Sigma_t^{1/2}
        \right)                        \\
         & =
        \frac2\lambda\sum_{t=1}^{\tau_{\mathrm{cov}}}
        \log\frac{\det(\overline\Sigma_{t+1}^{-1})}
        {\det(\overline\Sigma_t^{-1})} \\
         & =
        \frac2\lambda
        \log\frac{\det(\overline\Sigma_{\tau_{\mathrm{cov}}+1}^{-1})}
        {\det(\overline\Sigma_1^{-1})} \\
         & \leq
        \frac{2d}{\lambda}
        \log\left(
        \sigma^2
        \norm{\overline\Sigma_{\tau_{\mathrm{cov}}+1}^{-1}}_{\op}
        \right)                        \\
         & =
        \frac{4d}{\lambda}
        \log\left(
        \sigma
        \norm{\overline\Sigma_{\tau_{\mathrm{cov}}+1}^{-1}}_{\op}^{1/2}
        \right)                        \\
         & \leq
        \frac{4d}{\lambda}
        \log\left(1+n\right)\,,
    \end{align}
    where the last inequality
    uses \eqref{eq:predictable-spectrum},
    $\sigma\leq1$, and $\tau_{\mathrm{cov}}\leq n$.
    % Finally, by the lemma assumption and $\tau_{\mathrm{cov}}\leq n$ give
    % \begin{align}
    %     \log\left(1+\tau_{\mathrm{cov}}\lambda\eta L\sigma\right)
    %     \leq\log(1+n)\,,
    % \end{align}
    % which proves \eqref{eq:trace-budget}.
\end{proof}

\section{Feedback control}
\label{app:feedback}
\begin{proof}[Proof of \cref{lem:feedback}]
    Write
    \begin{align}
        r_t=\pi_\eps(X_t)=1+v(X_t),
        \qquad
        W_t=X_t-\mu_t.
    \end{align}
    \proofstepinline{1}
    Starting from the predictable second moment and conditioning first on
    $X_t$ gives
    \begin{align}
        \sum_{t=1}^{\tau_{\mathrm{cov}}}\E_t[Y_t^2]
         & =
        \sum_{t=1}^{\tau_{\mathrm{cov}}}
        \E_t\left[
            \E_t\left[
                \left(e(X_t)+r_t\xi_t\right)^2
                \,\middle|\,X_t
                \right]
            \right]
        \nonumber                        \\
         & \stackrel{\texttt{(a)}}{=}
        \sum_{t=1}^{\tau_{\mathrm{cov}}}
        \E_t\left[
            e(X_t)^2
            +
            r_t^2\E_t[\xi_t^2\mid X_t]
            \right]
        \nonumber                        \\
         & \stackrel{\texttt{(b)}}{\leq}
        \sum_{t=1}^{\tau_{\mathrm{cov}}}
        \E_t\left[
            e(X_t)^2+r_t^2
            \right]
        \nonumber                        \\
         & \stackrel{\texttt{(c)}}{\leq}
        4n
        +
        \frac{32}{\eps^2}
        \sum_{t=1}^{\tau_{\mathrm{cov}}}
        \E_t[v(X_t)^2].
        \label{eq:feedback-start}
    \end{align}
    Here \texttt{(a)} uses \eqref{eq:noise},
    and \texttt{(b)} follows from
    $\E_t[\xi_t^2\mid X_t]
        \leq
        \E_t[\exp(\xi_t^2)\mid X_t]-1
        \leq1$.
    For \texttt{(c)}, since $f\in[0,1]$ and $\eps\leq1/2$,
    \begin{align}
        0\leq e(X_t)
         & =
        r_tf(A_t)+\frac2\eps(r_t-1)
        \leq
        1+\frac3\eps v(X_t),
    \end{align}
    and therefore
    \begin{align}
        e(X_t)^2+r_t^2
         & \leq
        2+\frac{18}{\eps^2}v(X_t)^2
        +2+2v(X_t)^2
        \leq
        4+\frac{32}{\eps^2}v(X_t)^2.
    \end{align}

    Next, let $a=16\sqrt{\ell}$,
    then for every $t\leq\tau_{\mathrm{cov}}$,
    \begin{align}
        \E_t[v(X_t)^2]
         & =
        \E_t\left[
            v(X_t)^2\1\{v(X_t)\leq a\}
            \right]
        +
        \E_t\left[
            v(X_t)^2\1\{v(X_t)>a\}
            \right]
        \nonumber \\
         & \leq
        a\E_t[v(X_t)]
        +
        \E_t\left[
            v(X_t)^2\1\{v(X_t)>a\}
            \right].
        \label{eq:feedback-v2-split}
    \end{align}
    Since $\mu_t\in\K_\eps$, we have $v(\mu_t)=0$.  Convexity of $v$ and
    Gaussian integration by parts give
    \begin{align}
        \E_t[v(X_t)]
         & \stackrel{\texttt{(a)}}{\leq}
        \E_t\left[
            \ip{\nabla v(X_t)}{X_t-\mu_t}
            \right]
        \nonumber                        \\
         & \stackrel{\texttt{(b)}}{=}
        -\int_{\R^d}
        \ip{\nabla v(x)}
        {\Sigma_t\nabla\gamma_{\mu_t,\Sigma_t}(x)}
        \,dx
        \nonumber                        \\
         & =
        \sum_{i,j=1}^d
        (\Sigma_t)_{ij}
        \left\langle
        \partial_{ij}v,
        \gamma_{\mu_t,\Sigma_t}
        \right\rangle
        \nonumber                        \\
         & =
        \tr\left(
        \Sigma_t
        \left\langle
        D^2v,\gamma_{\mu_t,\Sigma_t}
        \right\rangle
        \right)
        \nonumber                        \\
         & \stackrel{\texttt{(c)}}{\leq}
        \eps\tr(\Sigma_tH_t).
        \label{eq:feedback-v-first-moment}
    \end{align}
    Here \texttt{(a)} is the supporting-hyperplane inequality at $X_t$,
    \texttt{(b)} uses
    $\Sigma_t\nabla\gamma_{\mu_t,\Sigma_t}(x)
        =-(x-\mu_t)\gamma_{\mu_t,\Sigma_t}(x)$
    and distributional integration by parts, and \texttt{(c)} uses
    $e(x) = \bar{e}(x) + \frac{1}{\eps} v(x)$ which, since $\bar{e}$ is
    convex and $v$ is convex, gives $D^2v \preceq \eps D^2e.$

    It remains to bound the tail term in
    \eqref{eq:feedback-v2-split},
    which we do by proving a sub-Gaussian tail for $v(X_t)$.
    By subadditivity of $\pi$, $\pi(\mu_t)\leq1-\eps$ since $\mu_t\in\K_\eps$, and $\eps\leq1/2$,
    \begin{align}
        v(X_t)
         & =
        \left(
        \frac{\pi(X_t)}{1-\eps}-1
        \right)_+
        \leq
        \left(
        \frac{\pi(X_t - \mu_t) + \pi(\mu_t)}{1-\eps}-1
        \right)_+
        \leq
        \frac{\pi(W_t)}{1-\eps}
        \leq
        2\pi(W_t).
        \label{eq:feedback-v-minkowski}
    \end{align}
    On the stopped interval,
    \eqref{eq:covariance-comparison} and monotonicity of
    $\overline\Sigma_t$ give
    $\Sigma_t
        \preceq
        2\overline\Sigma_t
        \preceq
        2\sigma^2I_d$.
    Therefore, since $\pi$ is $1$-Lipschitz, by Gaussian concentration we have
    \begin{align}
        \Pp_t\left(
        \pi(W_t)>
        \E_t[\pi(W_t)]+s
        \right)
        \leq
        \exp\left(-\frac{s^2}{2\|\Sigma_t\|_{\op}}\right)
        \leq
        \exp\left(-\frac{s^2}{4\sigma^2}\right).
        \label{eq:feedback-minkowski-tail}
    \end{align}
    Next, we bound the mean of $\pi(W_t)$.
    Let $Q \sim \mathcal{N}(0, 2\sigma^2 I_d - \Sigma_t)$
    so that $W_t + Q \sim \mathcal{N}(0, 2\sigma^2 I_d)$.
    Then, for $Z \sim \mathcal{N}(0, I_d)$ and $U\sim \text{Unif}(\mathbb{S}_1^d)$, we have
    \begin{align}
        \E_t[\pi(W_t)]
         & =
        \E_t[\pi(\E[W_t+Q\mid W_t])]
        \leq
        \E_t[\pi(W_t + Q)] \\
         & \leq
        \sqrt{2}\sigma
        \E_t[\pi(U)\cdot \|Z\|]
        \leq
        \sqrt{2d}\,\sigma\,\E[\pi(U)]
        \leq
        \sqrt{2d}\,\sigma M
        =1.
        \label{eq:feedback-minkowski-mean}
    \end{align}
    Since $M\geq d^{-1/2}$ and $\sigma M\sqrt{2d}=1$, we also have $\sigma^2\leq1/2$.
    Hence, for every $s\geq4$,
    \begin{align}
        \Pp_t\left(
        v(X_t)>s
        \right)
         & \leq
        \Pp_t\left(
        \pi(W_t)>\frac{s}{2}
        \right)
        \leq
        \exp\left(
        -\frac{(s/2-1)^2}{4\sigma^2}
        \right)
        \leq
        \exp\left(-\frac{s^2}{32}\right)\,,
        \label{eq:feedback-v-tail}
    \end{align}
    where the first inequality uses \eqref{eq:feedback-v-minkowski},
    the second uses \eqref{eq:feedback-minkowski-tail} and
    \eqref{eq:feedback-minkowski-mean}, and the last uses $\sigma^2\leq1/2$ and $s\geq4$.
    Since $a=16\sqrt\ell\geq4$, integrating this tail gives
    \begin{align}
        \E_t\left[
            v(X_t)^2\1\{v(X_t)>a\}
            \right]
         & =
        a^2\Pp_t(v(X_t)>a)
        +
        \int_a^\infty
        2s\Pp_t(v(X_t)>s)\,ds
        \nonumber \\
         & \leq
        a^2e^{-a^2/32}
        +
        \int_a^\infty2se^{-s^2/32}\,ds
        \nonumber \\
         & =
        (a^2+32)e^{-a^2/32}
        =
        (2^8\ell+32)e^{-8\ell}
        \leq
        e^{-\ell},
        \label{eq:feedback-v-tail-moment}
    \end{align}
    where the last inequality holds for $\ell\geq1$.
    Substituting
    \eqref{eq:feedback-v-first-moment} and
    \eqref{eq:feedback-v-tail-moment} into
    \eqref{eq:feedback-v2-split} gives
    \begin{align}
        \E_t[v(X_t)^2]
        \leq
        16\eps\sqrt\ell\,
        \tr(\Sigma_tH_t)
        +
        e^{-\ell}.
        \label{eq:feedback-v2}
    \end{align}
    We now return to \eqref{eq:feedback-start}, which can be further bounded as
    \begin{align}
        \sum_{t=1}^{\tau_{\mathrm{cov}}}\E_t[Y_t^2]
         & \stackrel{\texttt{(a)}}{\leq}
        4n
        +
        \frac{32}{\eps^2}
        \sum_{t=1}^{\tau_{\mathrm{cov}}}
        \E_t[v(X_t)^2]
        \\
         & \stackrel{\texttt{(b)}}{\leq}
        4n
        +
        \frac{2^9\sqrt\ell}{\eps}
        \sum_{t=1}^{\tau_{\mathrm{cov}}}
        \tr(H_t\Sigma_t)
        +
        \frac{32ne^{-\ell}}{\eps^2}
        \\
         & \stackrel{\texttt{(c)}}{\leq}
        4n
        +
        \frac{2^{12}d\sqrt\ell\log(1+n)}
        {\lambda\eta\eps}
        +
        \frac{32ne^{-\ell}}{\eps^2}
        \\
         & \stackrel{\texttt{(d)}}{\leq}
        4n
        +
        \frac{2^{12}n\ell}
        {c_\lambda c_\eta C_\eps \sqrt d\,M\,\ell^{3/2}}
        +
        n
        \stackrel{\texttt{(e)}}{\leq}
        6n
        \,,
        \label{eq:feedback-predictable-before-parameters}
    \end{align}
    Here \texttt{(a)} is \eqref{eq:feedback-start}, and
    \texttt{(b)} follows by applying \eqref{eq:feedback-v2} for each
    $t\leq\tau_{\mathrm{cov}}$ and using $\tau_{\mathrm{cov}}\leq n$.
    For \texttt{(c)}, the hypothesis of \cref{lem:trace} is satisfied by the
    assumption of the present lemma;  thus, \eqref{eq:covariance-comparison}
    and \eqref{eq:trace-budget} give
    \begin{align}
        \sum_{t=1}^{\tau_{\mathrm{cov}}}
        \tr(H_t\Sigma_t)
         & \leq
        2\sum_{t=1}^{\tau_{\mathrm{cov}}}
        \tr(H_t\overline\Sigma_t)
        \leq
        \frac{8d\log(1+n)}{\lambda\eta}.
        \label{eq:feedback-trace}
    \end{align}
    To justify \texttt{(d)} and \texttt{(e)}, the parameter definitions imply
    \begin{align}
        \frac{d\sqrt\ell\log(1+n)}
        {\lambda\eta\eps}
         & =
        \frac{
            n\log(1+n)
        }{
            c_\lambda c_\eta C_\eps
            \sqrt d\,M\,\ell^{3/2}
        }
        \leq
        2^{-12}
        n
        ,\quad\text{and}\quad
        \frac{e^{-\ell}}{\eps^2}
        =
        \frac{
            ne^{-\ell}
        }{
            C_\eps^2d^4M^2\ell^8
        }
        \leq1.
    \end{align}
    Indeed,
    $1<\log(1+n)\leq\ell$,
    $\sqrt d\,M\geq1$,
    $c_\lambda c_\eta C_\eps\geq 2^{12}$, and
    $ne^{-\ell}\leq1$; the fixed choice of the universal constants absorbs
    the remaining numerical factors.

    \proofstepinline{2}
    We next prove the realized bound.  Set
    \begin{align}
        B=\frac{32\ell}{\eps},
        \qquad
        \widetilde Y_t
        =
        Y_t\1\{|Y_t|\leq B\},
        \qquad
        \mathcal E_{\mathrm{tr}}
        =
        \left\{
        |Y_t|\leq B
        \text{ for every }t\leq\tau_{\mathrm{cov}}
        \right\}.
    \end{align}
    Since $Y_t=\widetilde Y_t$ on $\mathcal E_{\mathrm{tr}}$,
    \begin{align}
        \Pp\left(
        \sum_{t=1}^{\tau_{\mathrm{cov}}}Y_t^2
        >8n
        \right)
         & \leq
        \Pp(\mathcal E_{\mathrm{tr}}^c)
        +
        \Pp\left(
        \sum_{t=1}^{\tau_{\mathrm{cov}}}
        \widetilde Y_t^2
        >8n
        \right).
        \label{eq:feedback-realised-start}
    \end{align}

    First, we bound the probability of the truncation event $\mathcal E_{\mathrm{tr}}$.
    For $t\leq\tau_{\mathrm{cov}}$,
    \eqref{eq:feedback-minkowski-mean},
    \eqref{eq:feedback-minkowski-tail}, and
    $\sigma^2\leq1/2$ give
    $\Pp_t\left(
        \pi(W_t)>1+2\sqrt\ell
        \right)
        \leq e^{-2\ell}$,
    and the noise assumption gives
    $\Pp_t\left(
        |\xi_t|>\sqrt\ell
        \,\middle|X_t
        \right)
        \leq
        2e^{-\ell}$.
    If neither event occurs, then
    \eqref{eq:feedback-v-minkowski} yields
    \begin{align}
        r_t
        =1+v(X_t)
        \leq
        1+2\pi(W_t)
        \leq
        3+4\sqrt\ell
        \leq
        7\sqrt\ell,
    \end{align}
    and therefore
    \begin{align}
        |Y_t|
         & \leq
        e(X_t)+r_t|\xi_t|
        \leq
        r_t\left(\frac3\eps+|\xi_t|\right)
        \leq
        7\sqrt\ell
        \left(
        \frac3\eps+\sqrt\ell
        \right)
        \leq
        \frac{32\ell}{\eps}
        =B.
    \end{align}
    Since $\1\{t\leq\tau_{\mathrm{cov}}\}$ is
    $\F_{t-1}$-measurable, a union bound gives
    \begin{align}
        \Pp(\mathcal E_{\mathrm{tr}}^c)
         & \leq
        n\left(e^{-2\ell}+2e^{-\ell}\right)
        \leq
        3ne^{-\ell}
        \leq
        \frac{\delta}{16},
        \label{eq:feedback-truncation-event}
    \end{align}
    where the last inequality uses
    $\ell\geq\log(48n/\delta)$.

    Next, we bound the second term in \eqref{eq:feedback-realised-start} using a Freedman inequality.
    Define the martingale differences
    \begin{align}
        D_t
        =
        \1\{t\leq\tau_{\mathrm{cov}}\}
        \left(
        \widetilde Y_t^2
        -
        \E_t[\widetilde Y_t^2]
        \right).
    \end{align}
    Then
    $\E_t[D_t]
        =0$,
    $|D_t|\leq B^2$,
    and, using $\widetilde Y_t^4\leq B^2\widetilde Y_t^2$,
    \begin{align}
        \sum_{t=1}^n\E_t[D_t^2]
         & \leq
        \sum_{t=1}^n
        \1\{t\leq\tau_{\mathrm{cov}}\}
        \E_t[\widetilde Y_t^4]
        \leq
        B^2
        \sum_{t=1}^{\tau_{\mathrm{cov}}}
        \E_t[\widetilde Y_t^2]
        \leq
        B^2
        \sum_{t=1}^{\tau_{\mathrm{cov}}}
        \E_t[Y_t^2]
        \leq
        6nB^2.
        \label{eq:feedback-freedman-variance}
    \end{align}
    Applying the Freedman inequality \cref{lem:freedman} with
    \begin{align}
        b=B^2,
        \qquad
        s^2=6nB^2,
        \qquad
        z=\sqrt{2\log\frac{32}{\delta}},
    \end{align}
    and using $\log(32/\delta)\leq\ell$, we obtain, outside an event of
    probability at most $\delta/16$,
    \begin{align}
        \sum_{t=1}^nD_t
         & \leq
        \sqrt2\,s z+bz^2
        \leq
        2B\sqrt{6n\ell}
        +
        2B^2\ell.
        \label{eq:feedback-freedman}
    \end{align}
    Finally, the parameter definitions give
    $\frac{\ell^3}{\eps^2}
        =
        \frac{
            n
        }{
            C_\eps^2d^4M^2\ell^5
        }
        \leq n$.
    Hence, on event $\mathcal E_{\mathrm{tr}}$ and outside the Freedman event,
    \begin{align}
        \sum_{t=1}^{\tau_{\mathrm{cov}}}
        \widetilde Y_t^2
         & =
        \sum_{t=1}^n
        \1\{t\leq\tau_{\mathrm{cov}}\}
        \E_t[\widetilde Y_t^2]
        +
        \sum_{t=1}^nD_t
        \leq
        6n
        +
        2B\sqrt{6n\ell}
        +
        2B^2\ell
        \\
         &
        \leq
        6n
        +
        2^8\frac{\sqrt n \ell^{3/2}}{\eps}
        +
        2^{11}\frac{\ell^3}{\eps^2}
        \leq
        6n
        +
        \frac{2^8n}{C_\eps}
        +
        \frac{2^{11}n}{C_\eps}
        \leq
        8n\,,
        \label{eq:feedback-truncated-sum}
    \end{align}
    where the last inequality uses $C_\eps\geq2^{11}$.
    Combining
    \eqref{eq:feedback-realised-start},
    \eqref{eq:feedback-truncation-event}, and
    \eqref{eq:feedback-truncated-sum} gives
    \begin{align}
        \Pp\left(
        \sum_{t=1}^{\tau_{\mathrm{cov}}}Y_t^2
        >8n
        \right)
        \leq
        \frac{\delta}{16}
        +
        \frac{\delta}{16}
        =
        \frac{\delta}{8}.
    \end{align}
\end{proof}

\section{Concentration}
\label{app:concentration}
Throughout this section, write $\tau=\tau_{\mathrm{safe}}$.
The proof uses two scalar martingale inequalities.  The first one controls an
unbounded martingale by the sum of its realised and predictable quadratic
variations.

\begin{lemma}[Self-normalised martingale bound]
    \label{lem:self-normalised-martingale}
    Let $(D_t)_{t=1}^n$ be a square-integrable martingale-difference sequence, then for every $z,v>0$,
    \begin{align}
        \Pp\left(
        \left|\sum_{t=1}^n D_t\right|\geq z\quad\text{and}\quad
        \sum_{t=1}^n \left(D_t^2+\E_t[D_t^2]\right)\leq v
        \right)
        \leq
        2\exp\left(-\frac{z^2}{2v}\right).
        \label{eq:self-normalised-martingale}
    \end{align}
\end{lemma}

\begin{lemma}[Freedman's inequality]
    \label{lem:freedman}
    Let $(D_t)_{t=1}^n$ be a square-integrable
    martingale-difference sequence. Then, for every $b,s,z>0$,
    \begin{align}
        \Pp\left(
        \left|\sum_{t=1}^n D_t\right|
        \geq \sqrt{2}s z+bz^2,
        \quad
        \sum_{t=1}^n\E_t[D_t^2]\leq s^2,
        \quad
        \max_{1\leq t\leq n}|D_t|\leq b
        \right)
        \leq
        2\exp\left(-\frac{z^2}{2}\right).
    \end{align}
\end{lemma}

\begin{lemma}[Weighted Gaussian score moments]
    \label{lem:gaussian-score-localisation}
    Suppose that $t\leq\tau$, let $h_t$ be any
    $\F_{t-1}$-measurable unit vector, and put
    $G_t=\ip{h_t}{Z_t}$. Then
    \begin{align}
        \E_t\left[
            Y_t^2\left(
            \frac{G_t^2}{\ell}
            +\frac{G_t^4+1}{\ell^2}
            \right)
            \right]
         & \leq
        C\E_t[Y_t^2]+e^{-4\ell},
        \quad\text{and}\quad
        \E_t\left[
            Y_t^2(G_t^2-1)^2
            \right]
        \leq
        C\ell^2\E_t[Y_t^2]+e^{-4\ell}.
        \label{eq:H-score-second-moment}
    \end{align}
\end{lemma}

\begin{proof}
    Write $r_t=\pi_\eps(X_t)$. Since
    $Y_t=e(X_t)+r_t\xi_t$, conditional centring of $\xi_t$ gives
    \begin{align}
        \E_t[Y_t^2\mid X_t]
        =
        e(X_t)^2
        +r_t^2\E_t[\xi_t^2\mid X_t]
        \leq
        \frac{C}{\eps^2}r_t^2,
        \label{eq:Y2-conditional-simple}
    \end{align}
    where we used
    \begin{align}
        0\leq e(X_t)
        =r_tf(A_t)+\frac2\eps(r_t-1)
        \leq\frac{3r_t}{\eps}.
    \end{align}
    Moreover, by \eqref{eq:v-relation-to-minkowski},
    \begin{align}
        r_t
        =
        1+v(X_t)
        \leq
        1+2\pi(X_t-\mu_t).
    \end{align}
    Let $W_t=X_t-\mu_t$. Conditionally on $\F_{t-1}$,
    $W_t\sim\N(0,\Sigma_t)$, and, since $t\leq\tau$,
    $\Sigma_t\preceq2\sigma^2I_d$. Write
    $m_t=\E_t[\pi(W_t)]\leq1$. By
    \eqref{eq:gaussian-minkowski}, for every $s\geq1$,
    \begin{align}
        \Pp_t\bigl(\pi(W_t)>s\bigr)
        \leq
        \Pp_t\bigl(\pi(W_t)>m_t+s-1\bigr)
        \leq
        \exp\left(-\frac{(s-1)^2}{4\sigma^2}\right).
    \end{align}
    Integrating this tail and using $\sigma\leq1$ gives
    \begin{align}
        \E_t[\pi(W_t)^4]
         & =
        \int_0^\infty
        4s^3\Pp_t\bigl(\pi(W_t)>s\bigr)\,ds
        \leq
        1+
        \int_1^\infty
        4s^3
        \exp\left(-\frac{(s-1)^2}{4\sigma^2}\right)\,ds
        \leq C.
    \end{align}
    Consequently, using $(a+b)^4\leq8(a^4+b^4)$,
    \begin{align}
        \E_t[r_t^4]
        \leq
        8+128\E_t[\pi(W_t)^4]
        \leq C.
    \end{align}
    Let $\mathcal A_t=\{|G_t|\leq K\sqrt\ell\}$,
    where $K$ is a sufficiently large universal constant. On $\mathcal A_t$,
    \begin{align}
        \frac{G_t^2}{\ell}
        +\frac{G_t^4+1}{\ell^2}
        \leq C,
        \qquad
        (G_t^2-1)^2\leq C\ell^2.
        \label{eq:score-central-simple}
    \end{align}
    On the complement, since $G_t$ is determined by $X_t$ conditional on
    $\F_{t-1}$, \eqref{eq:Y2-conditional-simple} and Cauchy--Schwarz give
    \begin{align}
        \E_t\left[
            Y_t^2(1+G_t^4)\1\{\mathcal A_t^c\}
            \right]
         & \leq
        \frac{C}{\eps^2}
        \E_t\left[
            r_t^2(1+G_t^4)\1\{\mathcal A_t^c\}
            \right]
        \nonumber \\
         & \leq
        \frac{C}{\eps^2}
        \E_t[r_t^4]^{1/2}
        \E_t\left[
            (1+G_t^4)^2
            \1\{|G_t|>K\sqrt\ell\}
            \right]^{1/2}
        \nonumber \\
         & \leq
        \frac{C}{\eps^2}e^{-cK^2\ell}
        \leq e^{-4\ell}.
        \label{eq:score-tail-simple}
    \end{align}
    The penultimate inequality uses $G_t$'s conditional Gaussianity.  Indeed,
    another application of Cauchy--Schwarz and the Gaussian tail bound gives,
    with $a=K\sqrt\ell$,
    \begin{align}
         & \E_t\left[
            (1+G_t^4)^2\1\{|G_t|>a\}
            \right]^{1/2}
        \leq
        \E_t[(1+G_t^4)^4]^{1/4}
        \Pp_t(|G_t|>a)^{1/4}
        \leq
        Ce^{-a^2/8}
        =
        Ce^{-K^2\ell/8}.
    \end{align}
    For the last inequality in \eqref{eq:score-tail-simple}, the parameter
    definitions give
    \begin{align}
        \eps^{-1}
        =
        \frac{\sqrt n}{C_\eps d^2M\ell^4}
        \leq\sqrt n
        \leq e^{\ell/C_0}.
    \end{align}
    Thus it follows by choosing $K$ sufficiently large.
    Combining
    \eqref{eq:score-central-simple} and
    \eqref{eq:score-tail-simple} proves both claims.
\end{proof}

\begin{proof}[Proof of \cref{prop:concentration}]
    We first prove \eqref{eq:q-concentration}.
    To this end, define
    \begin{align}
        Q_t
        =\1\{t\leq\tau\}\widehat q_t(x_\star),
        \qquad
        D_t
        =Q_t-\E_t[Q_t]\,,
    \end{align}
    which is a martingale-difference sequence.
    The plan is to upper bound the quantity
    \begin{align}
        \sum_{t=1}^n
        \left(D_t^2 + \E_t[D_t^2]\right)
        \leq
        \sum_{t=1}^n
        \left(
        2Q_t^2 + 2\E_t[Q_t]^2 +
        \E_t[Q_t^2] - \E_t[Q_t]^2
        \right)
        \leq
        C
        \sum_{t=1}^n
        \left(
        Q_t^2 + \E_t[Q_t^2]
        \right)
        \,,
        \label{eq:D-to-Q}
    \end{align}
    and then apply \cref{lem:self-normalised-martingale} to bound $\left|\sum_{t=1}^\tau D_t\right| = \left|\sum_{t=1}^\tau (\widehat q_t(x_\star) - q_t(x_\star))\right|$.
    Note that $\E_t[Q_t] = \1\{t\leq\tau\}q_t(x_\star)$.

    Define $R_t=\normm{x_\star-\mu_t}{\Sigma_t^{-1}}$.
    Define the $\F_{t-1}$-measurable unit vector $h_t$ by
    \begin{align}
        h_t
        =
        \begin{cases}
            \frac{\Sigma_t^{-1/2}(x_\star-\mu_t)}{R_t},
             & R_t>0, \\
            e_1,
             & R_t=0,
        \end{cases}
    \end{align}
    Let $G_t=\ip{h_t}{Z_t}$, which is conditionally standard Gaussian.  The
    definitions of $\widehat g_t$ and $\widehat H_t$ give
    \begin{align}
        Q_t
        =Y_tL_t\1\{t\leq\tau\},
        \qquad\text{where}\quad
        L_t
        =R_tG_t+\frac{\lambda R_t^2}{2}(G_t^2-1).
        \label{eq:q-score-representation}
    \end{align}
    This identity remains valid when $R_t=0$, because both sides then vanish.
    A useful inequality is
    \begin{align}
        L_t^2 \1\{t\leq\tau\}
        \leq
        \frac{C}{\lambda^2}
        \left(
        \frac{G_t^2}{\ell}
        +\frac{G_t^4+1}{\ell^2}
        \right)
        \label{eq:q-score-square}
    \end{align}
    which follows from the focus condition \eqref{eq:concentration-focus}
    which ensures that $R_t\leq c/(\lambda\sqrt\ell)$ for $t\leq\tau$.
    We have
    \begin{align}
        \E_t[Q_t^2]
        \leq
        \E_t[Y_t^2L_t^2\1\{t\leq\tau\}]
        \stackrel{\texttt{(a)}}{\leq}
        \frac{C\1\{t\leq\tau\}}{\lambda^2}
        \E_t\left[Y_t^2 \left(
            \frac{G_t^2}{\ell}
            +\frac{G_t^4+1}{\ell^2}
            \right)\right]
        \stackrel{\texttt{(b)}}{\leq}
        \frac{C\1\{t\leq\tau\}}{\lambda^2}
        \left(\E_t[Y_t^2]+1\right),
        \label{eq:q-predictable-second-moment}
    \end{align}
    where \texttt{(a)} follows from \eqref{eq:q-score-square} and \texttt{(b)} follows from \cref{lem:gaussian-score-localisation}
    and the fact that $\ell\geq 1$.

    Let
    $\mathcal E_G
        =\left\{
        \max_{t\leq\tau}|G_t|\leq K\sqrt\ell
        \right\}$
    where $K$ is a sufficiently large universal constant.
    Since $G_t$ is conditionally standard Gaussian,
    \begin{align}
        \Pp(\mathcal E_G^c)
        \leq2ne^{-K^2\ell/2}
        \leq\frac{\delta}{64}.
        \label{eq:q-gaussian-maximum}
    \end{align}
    On $\mathcal E_G$, the focus condition \eqref{eq:concentration-focus} gives
    \begin{align}
        |L_t| \1\{t\leq \tau\} \leq |R_t||G_t| + \frac{\lambda R_t^2}{2}(|G_t^2| + 1)
        \leq
        \frac{C}{\lambda\sqrt\ell}K\sqrt\ell + \frac{\lambda}{2} \cdot \frac{C^2}{\lambda^2\ell}(K^2\ell + 1)
        \leq \frac{C}{\lambda}.
        \label{eq:q-score-bound}
    \end{align}
    Therefore, on $\mathcal E_G$,
    \begin{align}
        Q_t^2
        =
        Y_t^2L_t^2\1\{t\leq\tau\}
        \leq
        \frac{C\1\{t\leq\tau\}}{\lambda^2}
        Y_t^2\,.
        \label{eq:q-realised-score-energy}
    \end{align}
    Let $\mathcal E_Y$ be the event from \cref{lem:feedback} on which
    \begin{align}
        \sum_{t=1}^{\tau_{\mathrm{cov}}}Y_t^2\leq Cn\,,
    \end{align}
    then, since $\tau\leq\tau_{\mathrm{cov}}$, on
    $\mathcal E_Y\cap\mathcal E_G$,
    by \cref{lem:feedback} and \cref{eq:D-to-Q,eq:q-realised-score-energy,eq:q-predictable-second-moment} we have
    \begin{align}
        \sum_{t=1}^n
        \left(D_t^2+\E_t[D_t^2]\right)
         & \leq
        \frac{C}{\lambda^2}
        \sum_{t=1}^\tau
        \left(
        Y_t^2
        + \E_t[Y_t^2]
        +1
        \right)
        \leq\frac{Cn}{\lambda^2}
        \label{eq:q-total-self-normaliser}
    \end{align}
    Next, write
    \begin{align}
        \Pp\left(
        \left|\sum_{t=1}^nD_t\right|>
        \frac{C\sqrt{n\ell}}{\lambda}
        \right)
        \leq
        \Pp(\mathcal E_Y^c)
        +\Pp(\mathcal E_G^c)
        +
        \Pp\left(
        \left|\sum_{t=1}^nD_t\right|>
        \frac{C\sqrt{n\ell}}{\lambda},
        \ \sum_{t=1}^n(D_t^2+\E_t[D_t^2])
        \leq\frac{Cn}{\lambda^2}
        \right)
        \leq
        \frac{5\delta}{32}\,,
    \end{align}
    where we applied \cref{lem:self-normalised-martingale} with constant $C$ sufficiently large
    so that the last probability is at most $\delta/64$, which makes the total probability at most $\frac{\delta}{8} + \frac{\delta}{64} + \frac{\delta}{64} = \frac{5\delta}{32}$.
    This proves
    \eqref{eq:q-concentration}.

    Next we prove \eqref{eq:gradient-energy}.
    The definition of $\widehat g_t$ gives the exact identity
    $\normm{\widehat g_t}{\Sigma_t}^2=Y_t^2\norm{Z_t}^2$.
    Let
    $\mathcal E_Z
        =\left\{
        \max_{t\leq\tau}\norm{Z_t}^2\leq C(d+\ell)
        \right\}$.
    A chi-square tail bound and a union bound give
    \begin{align}
        \Pp(\mathcal E_Z^c)\leq\frac{\delta}{64}.
        \label{eq:Z-norm-union}
    \end{align}
    Hence, on $\mathcal E_Y\cap\mathcal E_Z$,
    \begin{align}
        \sum_{t=1}^{\tau}
        \normm{\widehat g_t}{\Sigma_t}^2
        \leq
        C(d+\ell)\sum_{t=1}^{\tau}Y_t^2
        \leq Cdn\ell,
    \end{align}
    which proves \eqref{eq:gradient-energy}.

    We now prove \eqref{eq:H-concentration}.
    We first record a deterministic range for the relevant metrics.
    The spectral estimate \eqref{eq:predictable-spectrum} gives
    \begin{align}
        \lambda_{\max}(\overline\Sigma_t^{-1})
        \leq
        \left(
        \sigma^{-1} + t
        \right)^2
        \leq
        C(d+n^2)
        \leq e^{C\ell}
        \qquad(t\leq\tau+1).
        \label{eq:concentration-spectrum}
    \end{align}
    Indeed, $B_1\subseteq\K$ implies $\pi(U)\leq1$, so
    $M\leq1$ and $\sigma^{-1}\leq\sqrt{2d}$.  Moreover,
    \eqref{eq:parameters-one}--\eqref{eq:parameters-two} and
    \eqref{eq:extension-lipschitz} give
    \begin{align}
        n\lambda\eta\Lip(e)
        \leq
        \frac{C c_\lambda c_\eta n}
        {C_\eps d^{5/2}M\ell^6}
        \leq Cn.
    \end{align}
    The last inequality in \eqref{eq:concentration-spectrum} follows from the
    definition of $\ell$.
    Let $\Gamma=2 e^{C\ell}$.
    Firstly, for every $t\leq\tau$, covariance comparison in \eqref{eq:covariance-comparison} gives
    \begin{align}
        I_d
        \preceq
        \frac{1}{2\sigma^2}I_d
        =
        \frac{1}{2} \overline\Sigma_1^{-1}
        \preceq
        \frac12\overline\Sigma_t^{-1}
        \preceq
        \Sigma_t^{-1}
        \preceq
        \frac32\overline\Sigma_t^{-1}
        \preceq
        \Gamma I_d\,,
        \label{eq:roundwise-spectral-range}
    \end{align}
    where we used \eqref{eq:concentration-spectrum}
    and $\sigma^{-2}=2dM^2\geq2$.
    Secondly, for $t=\tau+1$, the fact that
    $\overline\Sigma_{\tau+1}^{-1}\succeq\sigma^{-2}I_d$ and
    \eqref{eq:concentration-spectrum}
    give
    \begin{align}
        \kappa(\overline\Sigma_{\tau+1}^{-1})
        \leq \frac{
            e^{C\ell}
        }{
            \sigma^{-2}
        }
        \leq \Gamma\,.
    \end{align}
    Let $\mathcal U$ be a $(64\Gamma)^{-1}$-net of
    $\mathbb S^{d-1}$.
    By the standard volumetric bound for Euclidean covering numbers
    \cite[Corollary~4.2.13]{vershynin2018high}, it may be chosen so that
    \begin{align}
        \log|\mathcal U|
        \leq d\log(1+128\Gamma)
        \leq Cd\ell.
        \label{eq:H-net-cardinality}
    \end{align}
    Fix $u\in\mathcal U$ and define
    \begin{align}
        a_{t,u}
        =
        u^\top\Sigma_t^{-1}u,
        \qquad\text{and}\qquad
        G_{t,u}
        =
        \ip{
            \frac{
                \Sigma_t^{-1/2}u
            }{\sqrt{a_{t,u}}}
        }{Z_t}.
        \label{eq:H-standardisation}
    \end{align}
    Observe that $G_{t,u}$ is conditionally standard Gaussian.  Moreover,
    \begin{align}
        u^\top\widehat H_tu
        =
        Y_t\left(
        \ip{\Sigma_t^{-1/2}u}{Z_t}^2
        -
        u^\top\Sigma_t^{-1}u
        \right)
        =
        a_{t,u}Y_t(G_{t,u}^2-1).
        \label{eq:H-direction}
    \end{align}
    We first control
    $\sum_{t=1}^{\tau} Y_t^2(G_{t,u}^2-1)^2$
    for every $u\in\mathcal U$, outside an event of probability $\delta/64$.
    Let
    \begin{align}
        \mathcal C_t
        =
        \left\{
        |Y_t|\leq\frac{C\ell}{\eps},
        \quad
        \norm{Z_t}^2\leq C(d+\ell)
        \right\},
        \qquad
        W_{t,u}=Y_t^2(G_{t,u}^2-1)^2 \1\{\mathcal C_t\}
        \1\{t\leq\tau\}
        .
    \end{align}
    Since $G_{t,u}^2\leq\norm{Z_t}^2$, we have
    \begin{align}
        0\leq W_{t,u}\leq \Lambda,
        \qquad\text{where}\quad
        \Lambda=\frac{C\ell^2(d+\ell)^2}{\eps^2}\,.
        \label{eq:H-score-energy-increment}
    \end{align}
    Let $R_{t,u}=W_{t,u}-\E_t[W_{t,u}]$ which is a martingale difference sequence.
    Moreover,
    \begin{align}
        \sum_{t=1}^n
        \E_t[R_{t,u}^2]
        \leq
        \sum_{t=1}^n\E_t[W_{t,u}^2]
        \leq
        \Lambda\sum_{t=1}^n\E_t[W_{t,u}]
        =
        \Lambda\sum_{t=1}^{\tau}\E_t[W_{t,u}]
        \leq
        C\Lambda\ell^2n.
    \end{align}
    where the last inequality follows from
    \cref{lem:gaussian-score-localisation} and
    \eqref{eq:predictable-Y2} which give
    \begin{align}
        \sum_{t=1}^\tau
        \E_t[W_{t,u}]
        \leq
        \sum_{t=1}^\tau
        \E_t[Y_t^2(G_{t,u}^2-1)^2]
         & \leq
        \sum_{t=1}^{\tau}
        \left(
        C\ell^2\E_t[Y_t^2]+e^{-4\ell}
        \right)
        \leq C\ell^2n.
        \label{eq:H-score-energy-mean}
    \end{align}
    Now we are in a position to apply Freedman's inequality \cref{lem:freedman} to the martingale difference sequence $(R_{t,u})_{t=1}^n$ as
    \begin{align}
        |R_{t,u}|
        \leq \Lambda,         \qquad \text{and}\qquad
        \sum_{t=1}^n\E_t[R_{t,u}^2]
        \leq C\Lambda\ell^2n.
    \end{align}
    Let $J = \lceil \log_2(\Gamma) \rceil$
    and
    \begin{align}
        \alpha=\frac{\delta}{64|\mathcal U|}\,,
        \qquad\text{and}\qquad
        \rho
        =
        \log\frac{128|\mathcal U|(J+1)}{\delta}
        \leq Cd\ell.
    \end{align}
    For each fixed $u\in\mathcal U$, Freedman's inequality
    \cref{lem:freedman}, applied with $b=\Lambda$, and $s^2=C\Lambda\ell^2n$, gives, with probability at least $1-\alpha$,
    \begin{align}
        \left|
        \sum_{t=1}^nR_{t,u}
        \right|
         & \leq
        C\left(
        \sqrt{
            \Lambda\ell^2n
            \log\frac{2}{\alpha}
        }
        +
        \Lambda\log\frac{2}{\alpha}
        \right)
        \leq
        C\left(
        \sqrt{\Lambda\ell^2n\rho}
        +
        \Lambda\rho
        \right),
        \label{eq:H-score-energy-freedman}
    \end{align}
    where $\log\frac{2}{\alpha} = \log\frac{128|\mathcal U|}{\delta} \leq\rho$.
    A union bound over $u\in\mathcal U$ shows that
    \eqref{eq:H-score-energy-freedman} holds simultaneously for every
    $u\in\mathcal U$, except on an event of probability at most $|\mathcal U|\alpha=\frac{\delta}{64}$.
    Outside this exceptional event, for every $u\in\mathcal U$,
    \begin{align}
        \sum_{t=1}^nW_{t,u}
         & =
        \sum_{t=1}^n\E_t[W_{t,u}]
        +
        \sum_{t=1}^nR_{t,u} \\
         & \leq
        C\ell^2n
        +
        C\sqrt{\Lambda\ell^2n\rho}
        +
        C\Lambda\rho        \\
         & \leq
        C\ell^2n
        \left(
        1+
        \frac{(d+\ell)\sqrt\rho}{\eps\sqrt n}
        +
        \frac{(d+\ell)^2\rho}{\eps^2n}
        \right)             \\
         & \leq C\ell^2n\,,
        \label{eq:H-score-energy}
    \end{align}
    where the last inequality uses $d+\ell\leq2d\ell$, $\rho\leq Cd\ell$, and
    $\eps\sqrt n=C_\eps d^2M\ell^4$, which give
    \begin{align}
        \frac{(d+\ell)\sqrt\rho}{\eps\sqrt n}
        \leq
        \frac{C}{C_\eps\sqrt dM\ell^{5/2}}
        \leq\frac{C}{C_\eps}.
        \label{eq:H-score-energy-absorption}
    \end{align}
    On $\mathcal E_Y\cap\mathcal E_Z$, the event $\mathcal C_t$ holds
    for every $t\leq\tau$.
    Therefore, we have
    \begin{align}
        \sum_{t=1}^{\tau}
        Y_t^2(G_{t,u}^2-1)^2
        =
        \sum_{t=1}^{\tau}
        Y_t^2(G_{t,u}^2-1)^2
        \1\{\mathcal C_t\}
        =
        \sum_{t=1}^n
        W_{t,u}
        \leq C\ell^2n
        \label{eq:H-realised-score-energy}
    \end{align}
    simultaneously for every $u\in\mathcal U$.

    Define
    \begin{align}
        Q_{t,u}
        =
        \1\{t\leq\tau\}
        a_{t,u}Y_t(G_{t,u}^2-1), \qquad
        D_{t,u}
        =
        Q_{t,u}-\E_t[Q_{t,u}],
        \qquad\text{and}\qquad
        a_u^\star=\max_{t\leq\tau}a_{t,u}.
    \end{align}
    Let
    \begin{align}
        M_\tau
        =
        \sum_{t=1}^{\tau}(\widehat H_t-H_t).
    \end{align}
    Since $\E_t[\widehat H_t]=H_t$, we have
    $\sum_{t=1}^nD_{t,u} = u^\top M_\tau u$.
    By \eqref{eq:H-realised-score-energy},
    \begin{align}
        \sum_{t=1}^nQ_{t,u}^2
         & \leq
        (a_u^\star)^2
        \sum_{t=1}^{\tau}
        Y_t^2(G_{t,u}^2-1)^2
        \leq
        C(a_u^\star)^2\ell^2n.
        \label{eq:H-directional-realised}
    \end{align}
    Similarly, \cref{lem:gaussian-score-localisation} and
    \eqref{eq:predictable-Y2} give
    \begin{align}
        \sum_{t=1}^n\E_t[Q_{t,u}^2]
         & \leq
        (a_u^\star)^2
        \sum_{t=1}^{\tau}
        \left(
        C\ell^2\E_t[Y_t^2]+e^{-4\ell}
        \right)
        \leq
        C(a_u^\star)^2\ell^2n.
        \label{eq:H-directional-predictable}
    \end{align}
    Therefore, combining \eqref{eq:H-directional-realised} and
    \eqref{eq:H-directional-predictable} gives
    \begin{align}
        \sum_{t=1}^n
        \left(
        D_{t,u}^2+\E_t[D_{t,u}^2]
        \right)
        \leq
        \sum_{t=1}^n
        \left(
        Q_{t,u}^2+\E_t[Q_{t,u}^2]
        \right)
        \leq
        C(a_u^\star)^2\ell^2n.
        \label{eq:H-directional-self-normaliser}
    \end{align}
    Let $b_j=2^j$.
    For every $u\in\mathcal U$, choose the unique $j$ such that
    \begin{align}
        b_{j-1}<a_u^\star\leq b_j.
    \end{align}
    For each fixed pair
    $(u,j)$, \cref{lem:self-normalised-martingale} gives
    \begin{align}
         & \Pp\left(
        \left|\sum_{t=1}^nD_{t,u}\right|
        >Cb_j\ell\sqrt{n\rho},
        \quad
        \sum_{t=1}^n
        \left(
        D_{t,u}^2+\E_t[D_{t,u}^2]
        \right)
        \leq Cb_j^2\ell^2n
        \right)
        \leq
        2\exp\left(
        -\frac{C^2 b_j^2 \ell^2 n \rho}{C b_j^2 \ell^2 n}
        \right)
        \leq 2e^{-\rho}.
    \end{align}
    A union bound over $u\in\mathcal U$ and $0\leq j\leq J$ therefore has failure probability
    at most
    \begin{align}
        2|\mathcal U|(J+1)e^{-\rho}
        \leq \frac{\delta}{64}.
    \end{align}
    Therefore, outside an event of probability $\delta/64$,
    \begin{align}
        |u^\top M_\tau u|
        \leq
        C a_u^\star\ell\sqrt{n\rho}
        \leq
        \Xi
        \max_{t\leq\tau}u^\top\Sigma_t^{-1}u,
        \qquad\text{where}\quad
        \Xi =C\ell^{3/2}\sqrt{dn},
        \label{eq:H-net}
    \end{align}
    simultaneously for every $u\in\mathcal U$.

    It remains to extend \eqref{eq:H-net} from the finite net
    $\mathcal U$ to the whole sphere.
    Set $A=\frac32\overline\Sigma_{\tau+1}^{-1}$.
    For every $t\leq\tau$, covariance comparison and monotonicity of
    the predictable inverse covariance give
    \begin{align}
        \Sigma_t^{-1}
        \preceq
        \frac32\overline\Sigma_t^{-1}
        \preceq
        \frac32\overline\Sigma_{\tau+1}^{-1}
        =
        A.
        \label{eq:H-terminal-dominates}
    \end{align}
    Consequently, for every $u\in\mathcal U$,
    we have
    $\max_{t\leq\tau}u^\top\Sigma_t^{-1}u
        \leq
        u^\top Au$.
    Combining this with \eqref{eq:H-net} gives
    \begin{align}
        |u^\top M_\tau u|
        \leq
        \Xi
        \max_{t\leq\tau}u^\top\Sigma_t^{-1}u
        \leq
        \Xi u^\top Au
        \qquad(u\in\mathcal U).
        \label{eq:H-net-relative}
    \end{align}
    Moreover, scaling does not change the condition number, so we have
    $\kappa(A) = \kappa(\overline\Sigma_{\tau+1}^{-1}) \leq\Gamma$.
    Let $r=\frac{1}{64\Gamma}$.
    Fix $v\in\mathbb S^{d-1}$ and choose $u\in\mathcal U$ such that
    $\norm{u-v}\leq r$. Since $M_\tau$ is symmetric and
    $\norm{u}=\norm{v}=1$,
    \begin{align}
        \left|
        v^\top M_\tau v-u^\top M_\tau u
        \right|
        =
        \left|
        (v-u)^\top M_\tau v
        +
        u^\top M_\tau(v-u)
        \right|
        \leq
        2r\norm{M_\tau}_{\op}.
        \label{eq:H-net-approximation}
    \end{align}
    It follows from \eqref{eq:H-net-relative} that
    \begin{align}
        |v^\top M_\tau v|
        \leq
        \Xi u^\top Au
        +
        2r\norm{M_\tau}_{\op}
        \leq
        \Xi \lambda_{\max}(A)
        +
        2r\norm{M_\tau}_{\op}.
    \end{align}
    Taking the supremum over $v\in\mathbb S^{d-1}$ and using the
    symmetry of $M_\tau$ gives
    \begin{align}
        \norm{M_\tau}_{\op}
        \leq
        \Xi \lambda_{\max}(A)
        +
        2r\norm{M_\tau}_{\op}
        \leq
        2\Xi \lambda_{\max}(A),
        \label{eq:H-operator-bound}
    \end{align}
    where the last inequality uses $2r\leq1/2$.

    We now return to \eqref{eq:H-net-approximation}. Since
    $\norm{u-v}\leq r$,
    \begin{align}
        |u^\top Au-v^\top Av|
         & =
        \left|
        (u-v)^\top Au
        +
        v^\top A(u-v)
        \right|
        \leq
        2r\lambda_{\max}(A).
        \label{eq:H-A-net-approximation}
    \end{align}
    Therefore, by
    \cref{eq:H-net-relative,eq:H-operator-bound,eq:H-A-net-approximation},
    \begin{align}
        |v^\top M_\tau v|
         & \leq
        \Xi u^\top Au
        +
        2r\norm{M_\tau}_{\op}   \\
         & \leq
        \Xi v^\top Av
        +
        2\Xi r\lambda_{\max}(A)
        +
        4\Xi r\lambda_{\max}(A) \\
         & =
        \Xi v^\top Av
        +
        6\Xi r\lambda_{\max}(A).
    \end{align}
    Since $\kappa(A)\leq\Gamma$ and $\norm{v}=1$,
    \begin{align}
        \lambda_{\max}(A)
         & \leq
        \Gamma\lambda_{\min}(A)
        \leq
        \Gamma v^\top Av.
    \end{align}
    Hence
    \begin{align}
        6r\lambda_{\max}(A)
        \leq
        6r\Gamma v^\top Av
        =
        \frac{6}{64}v^\top Av
        \leq
        v^\top Av,
    \end{align}
    and consequently
    \begin{align}
        |v^\top M_\tau v|
        \leq
        2\Xi v^\top Av.
    \end{align}
    Since this holds for every $v\in\mathbb S^{d-1}$, we have
    \begin{align}
        -2\Xi A
        \preceq
        M_\tau
        \preceq
        2\Xi A.
    \end{align}
    Recalling that
    \begin{align}
        A=\frac32\overline\Sigma_{\tau+1}^{-1},
        \qquad
        \Xi =C\ell^{3/2}\sqrt{dn},
        \qquad
        M_\tau=\sum_{t=1}^{\tau}(\widehat H_t-H_t),
    \end{align}
    and absorbing numerical constants into $C$, we obtain
    \begin{align}
        -C\ell^{3/2}\sqrt{dn}\,
        \overline\Sigma_{\tau+1}^{-1}
        \preceq
        \sum_{t=1}^{\tau}(\widehat H_t-H_t)
        \preceq
        C\ell^{3/2}\sqrt{dn}\,
        \overline\Sigma_{\tau+1}^{-1}.
    \end{align}
    This proves \eqref{eq:H-concentration}.

    Finally, the total failure probability is at most
    \begin{align}
        \frac{\delta}{8}
        +
        5\frac{\delta}{64}
        =
        \frac{13\delta}{64}
        <
        \frac{\delta}{4}.
    \end{align}
    Here $\delta/8$ accounts for $\mathcal E_Y$, while the five terms
    of size $\delta/64$ account for $\mathcal E_G$, the
    self-normalised bound used for \eqref{eq:q-concentration},
    $\mathcal E_Z$, \eqref{eq:H-score-energy}, and the final
    self-normalised Hessian bound. This completes the proof.
\end{proof}





\section{Proof of \cref{prop:stability}}
\label{app:stability}

\begin{proof}[Proof of \cref{prop:stability}]
    We first verify the hypotheses of \cref{prop:bracket}.  Since
    $d^{3/2}M\geq d$, $\log(2e^\ell)\leq2\ell$, and \eqref{eq:concentration-focus}
    holds on the stopped interval,
    \begin{align}
        \lambda\left(d+\log(2e^\ell)\right)
         & \leq Cc_\lambda\leq c_0,
         &
        \lambda
        \normm{\mu_t-x_\star}{\Sigma_t^{-1}}
        \sqrt{\log(2e^\ell)}
         & \leq C\sqrt b\leq c_0.
    \end{align}
    Our choices also give
    $\lambda\leq c_\lambda\leq1/8$ and, after increasing $C_0$,
    $e^{-\ell}<1/2$.  Thus \cref{prop:bracket} applies with
    $\beta=e^{-\ell}$.

    Work on the event in \cref{prop:concentration}, which has probability at
    least $1-\delta/4\geq1-\delta/2$.

    We first show that the covariance condition does not trigger at
    $\tau_{\mathrm{safe}}$.  Since every earlier covariance comparison holds,
    starting from the precision at this round gives
    \begin{align}
        \Sigma_{\tau_{\mathrm{safe}}}^{-1}
        +\lambda\eta\widehat H_{\tau_{\mathrm{safe}}}
         & =
        \Sigma_1^{-1}
        +\lambda\eta\sum_{t=1}^{\tau_{\mathrm{safe}}}\widehat H_t
        =
        \overline\Sigma_{\tau_{\mathrm{safe}}+1}^{-1}
        +\lambda\eta\sum_{t=1}^{\tau_{\mathrm{safe}}}
        (\widehat H_t-H_t).
    \end{align}
    The parameter definitions give
    \begin{align}
        C\lambda\eta\ell^{3/2}\sqrt{dn}
        =
        \frac{Cc_\lambda c_\eta}{\sqrt\ell}
        \leq\frac14,
    \end{align}
    after taking $c_\lambda c_\eta$ sufficiently small.  Combining this with
    \eqref{eq:H-concentration} gives
    \begin{align}
        \frac34\overline\Sigma_{\tau_{\mathrm{safe}}+1}^{-1}
        \preceq
        \Sigma_{\tau_{\mathrm{safe}}}^{-1}
        +\lambda\eta\widehat H_{\tau_{\mathrm{safe}}}
        \preceq\frac54\overline\Sigma_{\tau_{\mathrm{safe}}+1}^{-1}.
        \label{eq:strict-comparison}
    \end{align}
    Thus the proposed precision at round $\tau_{\mathrm{safe}}$ is positive
    definite, so the algorithm performs the update and
    $\Sigma_{\tau_{\mathrm{safe}}+1}^{-1}
        =\Sigma_{\tau_{\mathrm{safe}}}^{-1}
        +\lambda\eta\widehat H_{\tau_{\mathrm{safe}}}$; moreover, the
    covariance comparison holds at this round.

    The covariance comparisons for $t<\tau_{\mathrm{safe}}$ and
    \eqref{eq:strict-comparison} at $t=\tau_{\mathrm{safe}}$ now give, for
    every $t\leq\tau_{\mathrm{safe}}$,
    \begin{align}
        \Sigma_{t+1}^{-1}
        \succeq
        \frac12\overline\Sigma_{t+1}^{-1}
        \succeq
        \frac12\overline\Sigma_t^{-1}
        \succeq
        \frac13\Sigma_t^{-1}.
    \end{align}
    Inverting this inequality and then applying \eqref{eq:gradient-energy}
    yields
    \begin{align}
        \sum_{t=1}^{\tau_{\mathrm{safe}}}
        \normm{\widehat g_t}{\Sigma_{t+1}}^2
        \leq
        3\sum_{t=1}^{\tau_{\mathrm{safe}}}
        \normm{\widehat g_t}{\Sigma_t}^2
        \leq Cdn\ell.
        \label{eq:next-gradient-energy}
    \end{align}
    Likewise, the lower comparison in
    \eqref{eq:stability-stopped-invariants} gives
    $\Sigma_t\preceq2\overline\Sigma_t$.  Since
    $\tau_{\mathrm{safe}}\leq\tau_{\mathrm{cov}}$ and
    $\tr(H_t\overline\Sigma_t)\geq0$, \cref{lem:trace} gives
    \begin{align}
        \eta\sum_{t=1}^{\tau_{\mathrm{safe}}}\tr(H_t\Sigma_t)
        \leq
        2\eta\sum_{t=1}^{\tau_{\mathrm{safe}}}
        \tr(H_t\overline\Sigma_t)
        \leq
        2\eta\sum_{t=1}^{\tau_{\mathrm{cov}}}
        \tr(H_t\overline\Sigma_t)
        \leq\frac{Cd\ell}{\lambda}.
        \label{eq:stability-stopped-trace}
    \end{align}

    We now bound the quadratic term appearing in the focus potential.  For
    every $t\leq\tau_{\mathrm{safe}}$,
    \cref{cor:regret-switch,prop:bracket} gives
    \begin{align}
        q_t(x_\star)
         & =
        -\bigl(\E_t[e(X_t)]-e(x_\star)\bigr)
        +\tr(H_t\Sigma_t)
        +\frac\lambda2\normm{\mu_t-x_\star}{H_t}^2
        -\E_t[D_e(x_\star\,\|\,X_t)]
        \nonumber \\
         & \leq
        -\bigl(\E_t[e(X_t)]-e(x_\star)\bigr)
        +C\tr(H_t\Sigma_t)
        +C\lambda\Lip(e)e^{-\ell}
        \norm{\mu_t-x_\star}
        \normm{\mu_t-x_\star}{\Sigma_t^{-1}}.
        \label{eq:q-upper}
    \end{align}
    Moreover, $\mu_t,x_\star\in\K_\eps\subseteq B_{2d}$ and
    \eqref{eq:stability-stopped-invariants} give
    \begin{align}
        \sum_{t=1}^{\tau_{\mathrm{safe}}}
        \lambda\Lip(e)e^{-\ell}
        \norm{\mu_t-x_\star}
        \normm{\mu_t-x_\star}{\Sigma_t^{-1}}
         & \leq
        4dn\lambda\Lip(e)e^{-\ell}\sqrt{2B}
        =
        \frac{Cdn\Lip(e)e^{-\ell}}{\sqrt\ell}
        \leq
        \frac{Cd n^{3/2}e^{-\ell}}{\sqrt\ell}
        \leq C.
        \label{eq:bracket-tail-sum}
    \end{align}
    The penultimate inequality uses \eqref{eq:extension-lipschitz} and
    \begin{align}
        \eps^{-1}
        =
        \frac{\sqrt n}{C_\eps d^2M\ell^4}
        \leq\sqrt n,
    \end{align}
    while the last follows from the definition of $\ell$ after increasing
    $C_0$.

    It remains to show that the focus condition does not trigger at
    $\tau_{\mathrm{safe}}$.  Starting from the focus potential at this round,
    summing \eqref{eq:ons-step}, and using $\widehat q_t(\mu_t)=0$ gives
    \begin{align}
        \frac12
        \normm{\mu_{\tau_{\mathrm{safe}}+1}-x_\star}
        {\Sigma_{\tau_{\mathrm{safe}}+1}^{-1}}^2
         & \leq
        \frac12\normm{\mu_1-x_\star}{\Sigma_1^{-1}}^2
        +\eta\sum_{t=1}^{\tau_{\mathrm{safe}}}\widehat q_t(x_\star)
        +\frac{\eta^2}{2}
        \sum_{t=1}^{\tau_{\mathrm{safe}}}
        \normm{\widehat g_t}{\Sigma_{t+1}}^2
        \\
         & =
        \frac12\normm{\mu_1-x_\star}{\Sigma_1^{-1}}^2
        +\eta\sum_{t=1}^{\tau_{\mathrm{safe}}}q_t(x_\star)
        +\eta\sum_{t=1}^{\tau_{\mathrm{safe}}}
        \bigl(\widehat q_t(x_\star)-q_t(x_\star)\bigr)
        +\frac{\eta^2}{2}
        \sum_{t=1}^{\tau_{\mathrm{safe}}}
        \normm{\widehat g_t}{\Sigma_{t+1}}^2
        \\
         & \leq
        \frac12\normm{\mu_1-x_\star}{\Sigma_1^{-1}}^2
        -\eta\sum_{t=1}^{\tau_{\mathrm{safe}}}
        \bigl(\E_t[e(X_t)]-e(x_\star)\bigr)
        +C\eta\sum_{t=1}^{\tau_{\mathrm{safe}}}\tr(H_t\Sigma_t)
        \nonumber \\
         & \qquad
        +C\eta\sum_{t=1}^{\tau_{\mathrm{safe}}}
        \lambda\Lip(e)e^{-\ell}
        \norm{\mu_t-x_\star}
        \normm{\mu_t-x_\star}{\Sigma_t^{-1}}
        \nonumber \\
         & \qquad
        +\eta\left|
        \sum_{t=1}^{\tau_{\mathrm{safe}}}
        \bigl(\widehat q_t(x_\star)-q_t(x_\star)\bigr)
        \right|
        +\frac{\eta^2}{2}
        \sum_{t=1}^{\tau_{\mathrm{safe}}}
        \normm{\widehat g_t}{\Sigma_{t+1}}^2
        \\
         & \leq
        \frac12\normm{\mu_1-x_\star}{\Sigma_1^{-1}}^2
        +C\eta^2dn\ell
        +\frac{Cd\ell}{\lambda}
        +\frac{C\eta\sqrt{n\ell}}{\lambda}
        +C\eta
        \label{eq:focus-budget}
    \end{align}
    The last inequality uses
    \eqref{eq:stationarity-sign}, \eqref{eq:stability-stopped-trace},
    \eqref{eq:bracket-tail-sum}, \eqref{eq:q-concentration}, and
    \eqref{eq:next-gradient-energy}.  Starting from the right-hand side of
    \eqref{eq:focus-budget} and substituting the parameters gives
    \begin{align}
         & \frac12\normm{\mu_1-x_\star}{\Sigma_1^{-1}}^2
        +C\eta^2dn\ell
        +\frac{Cd\ell}{\lambda}
        +\frac{C\eta\sqrt{n\ell}}{\lambda}
        +C\eta
        \\
         & \quad\leq
        C d^3M^2
        +C c_\eta^2d^3M^2\ell
        +\frac{C}{c_\lambda}d^{5/2}M\ell^3
        +\frac{Cc_\eta}{c_\lambda}d^{5/2}M^2\ell^{5/2}
        +\frac{Cc_\eta dM}{\sqrt n}
        \\
         & \quad\leq
        C\left(
        1+c_\eta^2+\frac1{c_\lambda}
        +\frac{c_\eta}{c_\lambda}+c_\eta
        \right)d^3M^2\ell^3
        \\
         & \quad\leq
        \frac{b}{c_\lambda^2}d^3M^2\ell^3
        =B.
    \end{align}
    The penultimate inequality uses $\sqrt dM\geq1$ and $\ell\geq1$.
    The last follows by first taking $c_\lambda$ sufficiently small relative
    to $b$ and then taking $c_\eta$ sufficiently small.
    All requirements are compatible by fixing the constants in the order
    $b$, $c_\lambda$, $c_\eta$, $C_\eps$, and finally $C_0$.
    Thus the focus condition also holds at round $\tau_{\mathrm{safe}}$.

    If either defining failure set were nonempty, its earliest element would
    be $\tau_{\mathrm{safe}}$.  The covariance and focus calculations show
    that neither condition fails there, so both failure sets are empty.
    Consequently,
    $\tau_{\mathrm{cov}}=\tau_{\mathrm{focus}}
        =\tau_{\mathrm{safe}}=n$ and
    both stopped invariants hold through time $n+1$.
    % The stopped concentration bound \eqref{eq:q-concentration} therefore gives
    % \eqref{eq:stability-q-budget}, while \eqref{eq:next-gradient-energy} and
    % \eqref{eq:stability-stopped-trace} give
    % \eqref{eq:stability-gradient-budget} and
    % \eqref{eq:stability-trace-budget}, respectively.
\end{proof}

\end{document}
